
\documentclass[11pt]{amsart}


\usepackage{palatino}
\usepackage{amsfonts}
\usepackage{amssymb}
\usepackage{amscd}
\usepackage{xypic}
\usepackage[breaklinks,bookmarksopen,bookmarksnumbered]{hyperref}
\usepackage{graphicx}
\usepackage{float}
\usepackage{comment}
\usepackage{xcolor}




\renewcommand{\baselinestretch}{1}
\renewcommand{\labelenumi}{(\roman{enumi})}
\newcommand{\dblq}{{/\!/}}
%\pagestyle{plain}
\setlength{\parindent}{.4 in}
\setlength{\textwidth}{6.2 in} \setlength{\topmargin} {-.3 in}
\setlength{\evensidemargin}{0 in}
%\theoremstyle{remark}
\newtheorem*{introtheorem}{Theorem}
\newtheorem*{introcorollary}{Corollary}
\newtheorem*{introconjecture}{Conjecture}
\newtheorem{theorem}{Theorem}[section]
\newtheorem{proposition}[theorem]{Proposition}
\newtheorem{corollary}[theorem]{Corollary}
\newtheorem{lemma}[theorem]{Lemma}
\theoremstyle{definition}
\newtheorem{definition}[theorem]{Definition}
\newtheorem{remark}[theorem]{Remark}
\newtheorem{example}[theorem]{Example}
\newtheorem{sketch}[theorem]{Sketch}
\newtheorem{conjecture}[theorem]{Conjecture}
\newtheorem{conjecture/question}[theorem]{Conjecture/Question}
\newtheorem{question}[theorem]{Question}
\newtheorem{remark/definition}[theorem]{Remark/Definition}
\newtheorem{terminology/notation}[theorem]{Terminology/Notation}
\newtheorem{assumption}[theorem]{Assumption}
\setlength{\oddsidemargin}{0 in} \setlength{\footskip}{.3 in}
\setlength{\headheight}{.3 in} \setlength{\textheight}{8.5 in}
\def\ZZ{{\mathbb Z}}
\def\GG{{\textbf G}}
\def\QQ{{\mathbb Q}}
\def\PP{{\textbf P}}
\def\FF{{\textbf F}}
\def\OO{\mathcal{O}}
\def\cN{\mathcal{N}}
\def\cD{\mathcal{D}}
\def\cB{\mathcal{B}}
\def\cA{\mathcal{A}}
\def\F{\mathcal{F}}
\def\P{\mathcal{P}}
\def\W{\mathcal{W}}
\def\E{\mathcal{E}}
\def\G{\mathcal{G}}
\def\X{\mathcal{X}}
\def\cS{\mathcal{S}}
\def\K{\mathcal{K}}
\def\L{\mathcal{L}}
\def\I{\mathcal{I}}
\def\J{\mathcal{J}}
\def\cM{\mathcal{M}}
\def\cR{\mathcal{R}}
\def\ra{\mathfrak{Rat}}
\def\rr{\overline{\mathcal{R}}}
\def\cZ{\mathcal{Z}}
\def\cU{\mathcal{U}}
\def\cC{\mathcal{C}}
\def\cH{\mathcal{H}}
\def\hh{\mathfrak{H}_{\mathrm{scr}}}
\def\pxi{P_{\xi}}
\def\ww{\overline{\mathcal{W}}}
\def\ff{\overline{\mathcal{F}}}
\def\mm{\overline{\mathcal{M}}}
\def\ss{\overline{\mathcal{S}}}
\def\kk{\overline{\mathcal{K}}}
\def\dd{\overline{\mathcal{D}}}
%\def\pm{\widetilde{\mathcal{M}}}
\def\ps{\widetilde{\mathcal{S}^{\mp}}}
\def\AUX{\mathfrak G^1_5(\widetilde{\cS}_8^{\mp}/\widetilde{\cM}_8)}
\def\thet{\overline{\Theta}_{\mathrm{null}}}
\DeclareMathOperator{\codim}{codim}
\DeclareMathOperator{\Hilb}{Hilb}
\DeclareMathOperator{\PGL}{PGL}
\DeclareMathOperator{\Ker}{Ker}
\DeclareMathOperator{\Supp}{Supp}
\DeclareMathOperator{\Sing}{Sing}
\DeclareMathOperator{\Aut}{Aut}
\DeclareMathOperator{\Ext}{Ext}
\DeclareMathOperator{\Hom}{Hom}
\DeclareMathOperator{\Pic}{Pic}
\DeclareMathOperator{\Sim}{Sym}
\DeclareMathOperator{\Span}{Span}
\DeclareMathOperator{\rank}{rank}
\DeclareFontFamily{OT1}{pzc}{}
\DeclareFontShape{OT1}{pzc}{m}{it}{<-> s * [1.10] pzcmi7t}{}
\DeclareMathAlphabet{\mathpzc}{OT1}{pzc}{m}{it}




\begin{document}
\title{Counterexamples to Green's Conjecture via $E_6$-covers}
\author[M. Aprodu]{Marian Aprodu}
\address{M. Aprodu: 
Faculty of Mathematics and Computer Science, 
\hfill \newline\texttt{}
 \indent University of Bucharest,  and Simion Stoilow Institute of Mathematics\hfill \newline\texttt{}
 \indent P.O. Box 1-764,
RO-014700 Bucharest, Romania
}
\email{{\tt marian.aprodu@imar.ro}}




\author[G. Farkas]{Gavril Farkas}

\address{G. Farkas: Humboldt-Universit\"at zu Berlin, Institut f\"ur Mathematik,  Unter den Linden 6
\hfill \newline\texttt{}
 \indent 10099 Berlin, Germany} \email{{\tt farkas@math.hu-berlin.de}}


\author[J. Park]{Jinhyung Park}
\address{J. Park: Department of Mathematical Sciences, KAIST, 91 Daehak-ro, Yuseong-gu, 
\hfill \newline\texttt{}
 \indent 34141 Daejeon, Republic of Korea
}
\email{{\tt parkjh13@kaist.ac.kr}}

 





\begin{abstract}
We show that curves of genus $10$ which are $W(E_6)$-covers of $\PP^1$ fail Green's Conjecture on syzygies of canonical curves. This provides a $9$-dimensional family of counterexamples to Green's Conjecture. This class of curves also provides counterexamples to the Green-Lazarsfeld Secant Conjecture. 
\end{abstract}

\maketitle




\section{introduction}

Formulated in 1984, Green's Conjecture \cite{Gr} has guided much of the research in syzygies and resolutions in the last decades. It predicts that the complexity of \emph{every} smooth algebraic curve $C$ of genus $g$ in the form of its Clifford index, which declares hyperelliptic curves to be the simplest curves of a given genus, can be recognized in a precise way in the equations of its canonical embedding
$\phi_{\omega_C}\colon C\hookrightarrow \PP^{g-1}=\PP H^0(C,\omega_C)^{\vee}$. Precisely, if $$\Gamma_C(\omega_C)=\bigoplus_{n\geq 0} H^0(C, \omega_C^{n})$$ is the canonical ring of $C$ viewed as a module over the symmetric algebra $S:=\mbox{Sym } H^0(C, \omega_C)$, we denote by $K_{p,q}(C, \omega_C)\cong \mbox{Tor}^p_S(\Gamma_C(\omega_C), \mathbb C)_{p+q}$ the Koszul cohomology group of $p$-syzygies of weight $q$ of the canonical curve. The Clifford index of the curve $C$ is defined as the quantity
$$
\mbox{Cliff}(C):=\mbox{min}\Bigl\{\mbox{deg}(L)-2h^0(C,L)+2: L\in \mbox{Pic}^d(C), h^0(C,L), h^1(C, L)\geq 2\Bigr\}
$$
Green's Conjecture \cite{Gr}, \cite{GL} predicts the following equivalence
\begin{equation}\label{eq:Green_Conj}
K_{p,2}(C,\omega_C)=0 \Longleftrightarrow p<\mbox{Cliff}(C). 
\end{equation}

\vskip 3pt

The statement $K_{p,2}(C, \omega_C)\neq 0$ if $p\geq \mbox{Cliff}(C)$ follows from the Green--Lazarsfeld \emph{Non-Vanishing Theorem} \cite[Appendix]{Gr}. The non-trivial implication is thus the vanishing $K_{p,2}(C,\omega_C)=0$, for  $p<\mathrm{Cliff}(C)$. Green's Conjecture has been proved for smooth curves on $K3$ surfaces of Picard number one in two landmark papers of Voisin \cite{V1},\cite{V2}. This implies Green's Conjecture for general curves of every genus. Coupled with work of Hirschowitz-Ramanan \cite{HR},  it follows from \cite{V2}  that Green's Conjecture holds for \emph{every} smooth curve of \emph{odd} genus $g$ and maximal gonality $\frac{g+3}{2}$. Using this fact, a sufficient Brill-Noether theoretic sufficient condition for a curve to satisfy Green's Conjecture was established in \cite{Ap}. Precisely, every curve $C$ of genus $g$ and gonality $k$ satisfying the following linear growth condition 
\begin{equation}\label{eq:aprodu1}
\mbox{dim } W^1_{g-k+2}(C)=g-2k+2
\end{equation}
verifies Green's Conjecture. In particular a general $k$-gonal curve of genus $g$ satisfies Green's Conjecture. It has also been shown that  curves  lying  on  \emph{arbitrary} $K3$ surfaces \cite{AF1}, or various covers satisfying generality assumptions \cite{AF3} all satisfy Green's Conjecture. Using work of Schreyer and collaborators \cite{Sch}, \cite{Sag}, Green's Conjecture is known to hold for all curves of genus $g<10$. Green's Conjecture has been recently showed to hold for other classes of curves, see for instance \cite{Mo}.

\vskip 4pt

In recent years several radically different proofs of Green's Conjecture for generic curves have been put forward, via Koszul modules and degeneration to sections of tangential varieties of rational normal curve \cite{AFPRW}, \cite{Pa}, \cite{RS}, or via a major simplification \cite{K} of Voisin's original proof for curves on $K3$ surfaces. It has also been established that Green's Conjecture holds for general curves of genus $g$ in positive characteristic at least $\frac{g-1}{2}$. 

\vskip 4pt


\noindent{\bf{The Secant Conjecture.}} The Green-Lazarsfeld Secant Conjecture \cite{GL} is a generalization of Green's Conjecture to the case of non-special bundles.  It predicts that if $L\in \mbox{Pic}^d(C)$ is a non-special line bundle on a curve $C$ of genus $g$ such that
\begin{equation}\label{ineq:gl}
d\geq 2g+p+1-\mbox{Cliff}(C),
\end{equation}
then $K_{p,2}(C,L)= 0$ if and only if $L$ is $(p+1)$-very ample, that is, every effective divisor of degree $p+2$ on $C$
imposes independent conditions on the linear system $|L|$. Via Riemann-Roch the Secant Conjecture can be reformulated as the following equivalence:
\begin{equation}\label{conj:GL}
K_{p,2}(C,L)\neq 0\Longleftrightarrow L-\omega_C\in C_{p+2}-C_{2g-d+p},
\end{equation}
where the right hand side denotes the corresponding difference variety in $\mbox{Pic}^{d-2g+2}(C).$
The Secant Conjecture has been established for a general pair $(C,L)$ in
\cite[Theorem 1.3]{FK}. In the case $\mbox{deg}(L)=2g+p+1-c$, there are complete results for \emph{arbitrary} curves when $c=1$, see \cite[Theorem 2]{GL} and when $c=2$, see  \cite{Ag}. Also the case when the difference variety appearing in (\ref{conj:GL}) is a divisor in $\mbox{Pic}^{d-2g+2}(C)$ has been essentially solved in \cite{F}. 

\vskip 4pt

As in the case of Green's Conjecture it remained an open problem whether the Secant Conjecture holds for arbitrary $C$ and $L$. In this paper we find counterexamples to both Green's Conjecture and the Secant Conjecture for certain curves of genus $10$. 

\vskip 5pt

The curves we consider in this paper  are degree $27$ covers of the projective line 
$$f\colon C\longrightarrow \PP^1$$
whose monodromy group $M(f)$ is the Weyl group $W(E_6)$ of the symmetries of the $27$ lines on a cubic surface. Thus each fibre $f^{-1}(t)$ inherits the structure of the $27$ lines on a fixed cubic surface and we assume furthermore that $f$ is branched at $12$ points $t_1, \ldots, t_{12}\in \PP^1$  such that at each point $t_i$ the ramification is given by a root of the $E_6$-lattice. The ramification profile of $f$ over each point $t_i$ is described by the partition $(2^6, 1^{15})$ and the points in $f^{-1}(t_i)$ can be identified with the lines on a one-nodal cubic surface. Each such a curve $C$ has genus $10$. We denote by $\mathcal{H}_{E_6}$ the Hurwitz space parametrizing such covers $\bigl[f\colon C\rightarrow \PP^1,\ t_1+\cdots+t_{12}\bigr]$. Kanev \cite{K2} showed that $\mathcal{H}_{E_6}$ is irreducible. Our main result is the following:

\begin{theorem}\label{thm:main}
The general $E_6$-cover $[f\colon C\rightarrow \PP^1]\in \cH_{E_6}$ corresponds to a smooth curve $C$ of genus $10$, Clifford index $\mathrm{Cliff}(C)=4$ and satisfying 
$$K_{3,2}(C, \omega_C)\neq 0.$$
In particular,  $C$ is a counterexample to Green's Conjecture.
\end{theorem}

Theorem \ref{thm:main} apparently contradicts both Theorem \cite[Theorem 1.3]{AF2} and \cite[p.184]{Eis}, which present Green's Conjecture as having been solved for all curves $C$ with $\mbox{Cliff}(C)\leq 4$.  However, both these claims are typos. For instance the argument in \cite{AF2} is obviously based on \cite[Theorem 4.8]{AF2}, which assumes $g\geq 12$, and indeed Green's Conjecture holds for all $6$-gonal curves of genus $g\geq 12$. 

\vskip 4pt

We moreover show that for a general curve as in Theorem \ref{thm:main},
$\mbox{dim } K_{3,2}(C, \omega_C)=1$. An $E_6$-curve $C$ like in Theorem \ref{thm:main} enjoys a number of remarkable properties. In particular, it is a curve of maximal gonality $\mbox{gon}(C)=6$, but has $\mbox{dim } W^1_6(C)=1$ and thus fails to satisfy the condition (\ref{eq:aprodu1}). 

\vskip 4pt

The $E_6$-covers of $\PP^1$ branched at $24$ (rather than $12$ points) points play a prominent role in the uniformization \cite{ADFIO} via the Prym-Tyurin map of the moduli space $\cA_6$ of principally polarized abelian varieties of dimension $6$. The Weyl group $W(E_6)$ appears also as the Galois group of the generically finite cover $\mathcal{R}_6\rightarrow \cA_5$, see \cite{Don}, \cite{DS}. The fact that $W(E_6)$ is also intricately related to Green's Conjecture was a genuine surprise for us.

\vskip 4pt

Generally, to an $E_6$-cover of $\PP^1$ branched over any number of points, one can associate a Prym-Tyurin variety.  As we shall explain, the Prym-Tyurin variety associated to a genus $10$ cover of $E_6$-type is in fact \emph{trivial}, and $10$ is the only genus for which this happens. This fact can be explained via a distinguished theta-characteristic on such a curve $C$ and the exceptional properties of it are ultimately responsible for the non-vanishing of the Koszul cohomology group $K_{3,2}(C, \omega_C)$. 

\vskip 4pt

We also present a concrete parametrization of the Hurwitz space $\cH_{E_6}$ and we show that every such curve $C\rightarrow \PP^1$ appears as the curve of lines in a special pencil of cubic surfaces 
$$X_{\lambda}=Z(\lambda F_0-F_1)\subseteq \PP^3, \ \mbox{ where } \lambda\in \PP^1, F_0, F_1\in H^0(\PP^3, \OO_{\PP^3}(3))$$
such that the base curve $B:=Z(F_0, F_1)\subseteq \PP^3$ is a rational curve with $10$ nodes. As explained in Theorem \ref{thm:curve_of_lines}, the curve $C$ can be regarded as the normalization of the curve $\mathrm{Tris}(B)$ of trisecant lines to $B$. We show in Theorem \ref{thm:parametrization} that the general $E_6$-cover of genus $10$ appears in this way. 

\vskip 4pt

Each $E_6$-curve $f\colon C\rightarrow \PP^1$ as above is endowed with a correspondence $\Sigma\subseteq C\times C$ consisting of pairs $(x,y)$ such that $x\neq y$ and  $f(x)=f(y)$ and the lines corresponding to the points $x$ and $y$ are incident. We denote by $S\subseteq C\times C$ the residual \emph{skew} correspondence corresponding to pairs $(x,y)$ such that  $f(x)=f(y)$ and $x$ and $y$ correspond to skew lines. If $\pi_1, \pi_2\colon C\times C\rightarrow C$ are the projections, for a point $x\in C$ we set $\Sigma_x:=\Sigma\cdot \pi_1^{-1}(x)$ and $S_x:=S\cdot \pi_1^{-1}(x)$. Since every line on a smooth cubic surface has $10$ incident and $16$ skew lines, it follows that $\mbox{deg}\bigl(\Sigma_x\bigr)=10$ and $\mbox{deg}\bigl(S_x\bigr)=16$.  It follows there exists a distinguished even theta-characteristic $\vartheta \in \mbox{Pic}^9(C)$ such 
$$\Sigma_x\equiv \vartheta(x), \ \mbox{ and } \ S_x\equiv \omega_C(-2x), \ \mbox{ for every }x\in C,$$
in particular $f^*(\OO_{\PP^1}(1))=\vartheta^{3}$. With this notation, we have the equivalence $(x,y)\in \Sigma$ if and only if $H^0(C,\vartheta(x-y))\neq 0$, that is, in the language of \cite{DK}, \cite{FI}, the incidence correspondence $\Sigma$ is equal to the \emph{Scorza correspondence} associated to the pair $(C,\vartheta)$. 

\vskip 4pt

The one-dimensional family $W^1_6(C)$ of minimal degree pencils has now a beautiful description in terms of pairs of roots of the lattice $E_6$ as follows. Given two skew lines on a cubic surface, it is well-known there are precisely \emph{five} lines incident to both of them. Translating this fact into the language of $E_6$-covers, it follows that for a pair $(x,y)\in S$, the degree $10$ divisors $\Sigma_x$ and $\Sigma_y$ have $5$ common points in their support, that is, we can write 
$$\Sigma_x=D_{xy}+R_{xy} \ \ \mbox{ and } \ \ \Sigma_y=D_{xy}+R'_{xy},$$
where $R_{xy}$ and $R_{xy}'$ are disjoint and  $\mbox{deg}(R_{xy})=\mbox{deg}(R'_{xy})=5$. It follows that the degree $6$ divisors $R_{xy}+y$ and $R_{xy}'+x$ are linearly equivalent, thus they generate a pencil 
\begin{equation}\label{eq:pencils1}
A_{xy}:=\OO_C(R_{xy}+y)=\OO_C(R'_{xy}+x)\in W^1_6(C).
\end{equation}
We then show that $W^1_6(C)$ is parametrized by the degree $36$ cover $C_{36}\rightarrow \PP^1$ associated to $f\colon C\rightarrow \PP^1$ and parametrizing the \emph{double-sixes} in the fibres of $f$.

\vskip 5pt

\noindent{\bf{Scorza syzygies.}} To the theta-characteristic $\vartheta$ we can associate a smooth quadric $Q=Q_{\vartheta}\in \mbox{Sym}^2 H^0(C,\omega_C)$  containing the canonical curve $C\subseteq \PP^9$ and having the following remarkable property for pairs $(x,y)\in C\times C$ 
\begin{equation}\label{eq:polar_quad}
Q(x,y)=0 \Longleftrightarrow x=y \ \mbox{ or } (x,y)\in S.
\end{equation}
In other words, $Q$ detects the incidence of the lines corresponding to points $x,y$ lying in the same fibre of $f$. To define $Q$ we use that $\OO_{C\times C}(\Sigma)\cong \vartheta\boxtimes \vartheta(\Delta)$, where $\Delta\subseteq C\times C$ is the diagonal. Denoting by $\mathrm{sz}_{\vartheta}(x,y)$ the \emph{Szeg\H{o} kernel} of $\vartheta$, that is, the section of $\vartheta\otimes \vartheta(\Delta)$ whose zero locus is the curve $\Sigma$, if $(s_0, s_1)$ is a basis of $f^*H^0(\PP^1, \OO_{\PP^1}(1))$, 
we set 
\begin{equation}\label{eq:Scorza_quadric}
Q(x,y)=\frac{s_0(x)s_1(y)-s_0(y)s_1(x)}{\mathrm{sz}_{\vartheta}(x,y)}\in H^0\bigl(C\times C, (\omega_C\boxtimes \omega_C)(-\Delta)\bigr)^{S_2}\cong I_2(\omega_C).
\end{equation}
We refer to $Q$ as the Scorza quadric of the $E_6$-curve $f\colon C\rightarrow \PP^1$. To $Q$ we associate its quadratic rank $8$ vector bundle $\cH_Q=\Omega_{Q}^1|_C\otimes\omega_C$ on $C$, sitting in an exact sequence 
\begin{equation}\label{eq:quadrbunde}
0\longrightarrow \omega_C^{\vee}\longrightarrow M_{\omega_C}\longrightarrow \cH_Q\longrightarrow 0,
\end{equation}
where as usual $M_{\omega_C}:=\mbox{Ker}\bigl\{H^0(C,\omega_C)\otimes \OO_C\rightarrow \omega_C\bigr\}$ is the syzygy bundle. We refer to $\cH_Q$ as being \emph{Scorza vector bundle} associated to the $E_6$-curve $f\colon C\rightarrow \PP^1$. 

\vskip 4pt

Taking exterior powers in  (\ref{eq:quadrbunde}), we obtain an inclusion
$ \bigwedge^4 \cH_Q\hookrightarrow \bigwedge^5 M_{\omega_C}\otimes \omega_C$. The smooth Scorza quadric $Q$ endows $\mathcal H_Q$ with a nondegenerate symmetric complex-bilinear form.  The bundle $\mathcal H_Q$ admits a natural orthogonal structure. A normalized orientation defines the algebraic Hodge-star involution on its middle exterior power. We denote its two rank-$35$ eigensubbundles by $\mathcal F_Q^\pm$ such that 
\[
 \bigwedge^4\mathcal H_Q=\mathcal F_Q^+\oplus\mathcal F_Q^-.
\]


\vskip 4pt

Proposition \ref{prop:star} shows that the natural inclusion
\[
 \bigwedge^4\mathcal H_Q\longrightarrow
 \bigwedge^5 M_{\omega_C}\otimes\omega_C
\]
induces injective linear maps
\[
 H^0(C,\mathcal F_Q^+)\hookrightarrow K_{5,1}(C,\omega_C)\mbox{ and } H^0(C,\mathcal F_Q^-)\hookrightarrow K_{5,1}(C,\omega_C).
\]

The assertion follows from the kernel-bundle description of Koszul cohomology. Sections of $\bigwedge^4\mathcal H_Q$ give fifth linear cycles, while the Hodge-eigenvector condition excludes Koszul boundaries. Consequently, by canonical Koszul duality, the non--vanishing of $H^0(C,\bigwedge^4\mathcal H_Q)$ implies $K_{3,2}(C,\omega_C)\ne0$. The proof of our main theorem is therefore reduced to producing a nonzero section of $\bigwedge^4\mathcal H_Q$. We do so by using $16$ distinguished isotropic $4$-planes, see Proposition \ref{prop:6=16},
$$P_{x,y}\subseteq \cH_{Q|x},\  \mbox{ where } \ (x,y)\in S. $$
These $16$ four-planes are governed in a strict way by the Weyl group $W(D_5)$ describing the symmetries of the $16$ lines on a del Pezzo quartic surface. Due to this structure, it turns out that for every $x\in C$ there is a $3$-dimensional space of alternating forms on $\bigwedge^2 \cH_{Q|x}^{\vee}$ vanishing on these planes. This induces a rank $3$ subbundle 
$$\mathcal N\subseteq \bigwedge^2\mathcal H_Q^\vee.$$ 
The second symmetric power $\operatorname{Sym}^2\mathcal{N}$  naturally embeds into $\bigwedge^4\mathcal H_Q$. Then the restriction of the quadratic form to $\mathcal{N}$ provides us with the desired non--zero section of $\bigwedge^4\mathcal H_Q$.

\vskip 4pt

At the end of the paper, we discuss in Theorem \ref{th:def_GL} a counterexample to the Green-Lazarsfeld Secant Conjecture. It turns out that any $E_6$-cover $f\colon C\rightarrow \PP^1$ of genus $10$ possesses a line bundle $L\in \mbox{Pic}^{20}(C)$ which is $4$-very ample, yet $K_{3,2}(C,L)\neq 0$.


\vskip 9pt

\noindent {\bf Acknowledgments:} Farkas was  partially supported by the Berlin Mathematics Research Center MATH+, by the Einstein Foundation Berlin and by the ERC Advanced Grant SYZYGY (grant No. 834172). Aprodu is grateful to KAIST Daejeon for hospitality during the first stage of this work.




\section{The Hurwitz space of $E_6$-covers of genus $10$}

\subsection{Basic notation on syzygies.} We begin by recalling the usual description of the Koszul cohomology groups of curves. We choose a smooth curve $C$ and a  globally generated line bundle $L\in \mbox{Pic}^d(C)$ and consider the Koszul complex
\begin{align*}
\bigwedge^{p+1} H^0(C,L)\otimes H^0\bigl(C, L^{q-1}\bigr)\stackrel{d_{p+1,q-1}}\longrightarrow \bigwedge^p H^0(C,L)\otimes H^0\bigl(C, L^q\bigr) \\
\stackrel{d_{p,q}}\longrightarrow \bigwedge^{p-1} H^0(C,L)\otimes H^0(C,  L^{q+1}\bigr),
\end{align*}
where $d_{p,q}$ is the Koszul differential. Then $K_{p,q}(C,L):=\mbox{Ker}(d_{p,q})/\mbox{Im}(d_{p+1,q-1})$. With the help of the \emph{syzygy bundle}
$M_L:=\mbox{Ker}\bigl\{H^0(C,L)\otimes \mathcal{O}_C\rightarrow L\bigr\}$,
one has the following  interpretation of the Koszul cohomology groups
\begin{equation}\label{eq:kernel_bundle}
 K_{p,q}(C, L) \cong \frac{H^0\bigl(C, \bigwedge^p M_L\otimes L^q \bigr)}{\mathrm{Im}\big(d_{p+1,q-1}\big)}.
\end{equation}
When $L=\omega_C$, one has moreover the duality $K_{p,2}(C,\omega_C)\cong K_{g-p-2,1}(C, \omega_C)^{\vee}$.

\vskip 4pt
For instance if $C$ is a curve of genus $10$ and of maximal Clifford index $\mbox{Cliff}(C)=4$, Green's Conjecture for $C$ is equivalent to the vanishing 
$$K_{3,2}(C,\omega_C)=0\Longleftrightarrow K_{5,1}(C,\omega_C)=0.$$

\subsection{$E_6$-covers of the projective line.} We now  discuss basics on curves expressed as covers of $\PP^1$ with monodromy group $W(E_6)$.  Our standard reference for facts related to cubic surface and the $E_6$ lattice is \cite[Chapters 8,9]{Dol}.
We fix once and for all a smooth cubic surface $X$ and identify the Picard group $\mbox{Pic}(X)$  with the Lorentzian lattice $I^{1,6}$ endowed with the intersection product on $X$. We denote by $h\in \mbox{Pic}(X)$ the hyperplane class and by $a_1, \ldots, a_6$ the exceptional divisors corresponding to the six points one has to blow-up $\PP^2$ in order  to get $X$. Then $E_6:=K_X^{\perp}$,  where $K_X=-3h+a_1+\cdots+a_6$. 

\vskip 4pt

A vector  $r\in E_6$ with $r^2=-2$ is called a root. There are $36$ pairs of roots corresponding to
$\alpha_{ij}=a_i-a_j$, $\alpha_{ijk}=h-a_i-a_j-a_k$ and
$\alpha_{\max}=2h -a_1-\cdots -a_6$. Obviously, if $r\in E_6$ is a root then $-r$ is a root as well.  The Weyl group $W(E_6)$ is defined as the subgroup of the orthogonal group $O(E_6)$  generated by the reflections in the
roots. It has 51840 elements and $25$ irreducible representations.

\vskip 4pt

Lines on the cubic surface $X$ can  be identified with the exceptional vectors $\ell\in \mbox{Pic}(X)$ such that
$\ell^2=K_X\cdot \ell=-1$. As in \cite{Dol}, we use the following notation for the lines:
\begin{align*}
  &a_i,\ b_i = 2h-a_1-\dots-a_6+a_i, \ \mbox{ for } i=1, \ldots, 6 \mbox{ and } 
  c_{ij}= h-a_i-a_j, \  \mbox{ for } \quad 1\le i<j \le 6.
\end{align*}

For each root $r\in E_6$, there are $15$ exceptional vectors that are
orthogonal to it, $6$ exceptional vectors with $r\cdot \ell=1$ and $6$ vectors with
$r\cdot \ell=-1$. The collections of the $6$ pairs of exceptional vectors
non-orthogonal to a root vector are called \emph{double-sixes}. The
elements in each pair are exchanged by the reflection $w_r\in W(E_6)$ in the
root $r$. There are $36$ double-sixes, one for each pair $\pm r$ of roots. For example, the double-six for the root $r=\alpha_{\max}$ is
$\{a_1,a_2,\dotsc,a_6\}$, $\{b_1,b_2,\dotsc,b_6\}$. Since $W(E_6)$ acts transitively on the set of lines, this gives rise to an embedding $W(E_6)\subset S_{27}$ under which
 each reflection corresponds to a product of 6
transpositions. For example, the reflection in the maximal root
$r=\alpha_{\max}$ is the permutation $(a_1, b_1)\cdots (a_6, b_6)\in S_{27}$.

\begin{definition}\label{def:hurw}
Let $\cH_{E_6}$ be the Hurwitz space parametrizing covers $\bigl[f\colon C\rightarrow \PP^1, t_1, \ldots, t_{12}]$ having monodromy groups $M(f)=W(E_6)$ and branched over the points $t_i\in \PP^1$ such that the local monodromy around $t_i$ is given by a reflection $w_i$ corresponding to a root of $E_6$.
\end{definition}
One has two morphisms retaining the source and the branch divisor of each $E_6$-cover 
\begin{equation}\label{eq:br+m}
\xymatrix{
  & \cH_{E_6}
   \ar[dl]_{m} \ar[dr]^{\mathfrak{br}} & \\
   \cM_{10} & & \cM_{0,12}       \\
                 }
\end{equation}
The map $\mathfrak{br}$ is finite, hence $\mbox{dim}(\cH_{E_6})=9$. Using Kanev's results \cite{K2}, it also follows that $\cH_{E_6}$ is irreducible.

\vskip 4pt

We pick an $E_6$-cover $f\colon C\rightarrow \PP^1$ branched over the points $t_1, \ldots, t_{12}\in \PP^1$. Each fibre of $f$ is consistently identified  with the set of $27$ lines on our reference cubic surface $X$. The incidence of lines on $X$ on unramified fibres induces by passing to the closure a correspondence $\Sigma\subseteq C\times C$ of valence (degree) $10$. This correspondence is symmetric and disjoint from the diagonal $\Delta\subseteq C\times C$. We denote by $S\subseteq C\times C$ the residual correspondence parametrizing along each unramified fibre of $f$ skew lines in $X$. We have the following decomposition of the fibre product
$$C_{\times_{\PP^1}} C=\Sigma+S+\Delta.$$
For each point $x\in C$, we denote by $\Sigma_x$ respectively by $S_x$ the fibre of the respective correspondence over $x$. The correspondence $\Sigma$ induces an endomorphism $\sigma\colon \mbox{Pic}^0(C)\rightarrow  \mbox{Pic}^0(C)$ defined by the assignment $x\mapsto \Sigma_x\in \mbox{Div}^{10}(C)$ at the level of divisors. Using \cite{K1}, or \cite[\S1]{ADFIO}, this endomorphism  satisfies Kanev's quadratic relation 
\begin{equation}\label{eq:kanev}
(\sigma-1)(\sigma+5)=0\in \mbox{End}\bigl(\mbox{Pic}^0(C)\bigr).
\end{equation}
Crucially for our purposes, this endomorphism is trivial in genus $10$, that is, the associated Prym-Tyurin variety $PT(C, \Sigma)=\mbox{Im}(\sigma-\mathrm{id}_{\mathrm{Pic}^0(C)})\subseteq \mbox{Pic}^0(C)$ is trivial.

\begin{proposition}\label{prop:corresp1}
Keeping the notation as above, the following hold for an $E_6$-cover $f\colon C\rightarrow \PP^1$:
\begin{enumerate}
    \item $\sigma=\mathrm{id}_{\mathrm{Pic}^0(C)}.$
    \item The cover $f$ induces a theta-characteristic $\vartheta\in \mathrm{Pic}^9(C)$ with $H^0(C, \vartheta)=0$, such that  $\OO_C(\Sigma_x)\cong \vartheta(x)$,  for every $x\in C$. Furthermore, $\OO_C(S_x)\cong \omega_C(-2x)$, for every $x\in C$.
    \item $f^*(\OO_{\PP^1}(1))\cong \vartheta^{\otimes 3}$.
\end{enumerate}
\end{proposition}
\begin{proof}
Using Kanev's identity (\ref{eq:kanev}), one has the following decomposition into eigenspaces 
$$H^0(C,\omega_C)=H^0(C,\omega_C)^{(+1)}\oplus H^0(C,\omega_C)^{(-5)}.$$ 
We observe that $\Sigma$ is disjoint from $\Delta$. This is obvious in a fibre of $f$ that is not ramified and it only has to be explained over one of the ramification points. Assuming, without loss of generality, that $\alpha_{\mathrm{max}}$ is the root of $E_6$ giving the local monodromy of $f$ around this point. The associated double six interchanges the lines $a_i$ and $b_i$ for $i=1, \ldots, 6$, while leaving the lines $c_{ij}$ fixed. Notice, that the lines $a_i$ and $b_i$ which are identified for $i=1, \ldots, 6$ are all pairwise skew. Therefore $\Sigma\cap \Delta=\emptyset$.
\vskip 3pt

From the Lefschetz fixed point formula, we obtain that 
$$0=\Sigma\cdot \Delta=2\mbox{deg}(\Sigma)-2\mbox{tr}\bigl\{\sigma\colon H^0(C,\omega_C)\longrightarrow H^0(C,\omega_C)\bigr\},$$
see also \cite[1.10]{ADFIO}, from which we compute 
$$\mbox{dim } H^0(C,\omega_C)^{(-5)}=\frac{1}{6}\bigl(g(C)-\mbox{deg}(\Sigma)\bigr)=\frac{1}{6}\bigl(10-10\bigr)=0,$$
therefore $H^0(C,\omega_C)=H^0(C,\omega_C)^{(+1)}$ and the endomorphism $\sigma\in \mbox{End}\bigl(\mbox{Pic}^0(C)\bigr)$ s trivial. It follows that there exists a line bundle $\vartheta\in \mbox{Pic}^9(C)$ such that 
$$\OO_{C\times C}(\Sigma)\cong \vartheta\boxtimes \vartheta(\Delta).$$
Since $\Sigma$ and $\Delta$ are disjoint, by restricting this relation to $\Delta$, we obtain that $\vartheta^{2}\cong \omega_C$, that is, $\vartheta$ is a theta-characteristic and for every $x\in C$, we have that $\OO_C(\Sigma_x)\cong \vartheta(x)$. In particular, $\Sigma_x$ is an  element of the linear system $|\vartheta(x)|$ which does not contain the point $x$, which implies that $H^0(C,\vartheta)=0$. Indeed, assuming that $h^0(C,\vartheta)=r\geq 1$, then for a general point $x\in C$, we have that 
$h^0(C,\vartheta(-x))=r-1$, hence by Riemann-Roch $h^0(C,\vartheta(x))=r$. It follows then for dimension reasons that the injective map $H^0(C,\vartheta)\hookrightarrow H^0(C,\vartheta(x))$ is an isomorphism, in particular $x\in \mathrm{Bs}\ |\vartheta(x)|$, which is a contradiction. Therefore, $\vartheta$ is a non-effective spin structure on $C$.


\vskip 4pt

To prove (iii), we use that if $x,z,y\in f^{-1}(t)$ are points corresponding to a tritangent triangle on the cubic surface $X$, that is, the lines corresponding $x,y,z$ mutually intersect, then  
\begin{equation}\label{eq:tritangent1}
\Sigma_x+\Sigma_y+\Sigma_z=x+y+z+f^{-1}(t)\in \mbox{Div}^{30}(C).
\end{equation}
Indeed, since the action of $W(E_6)$ on the set of $45$ tritangent planes is transitive, without loss of generality, we may assume $x=a_1$, $y=b_2$ and $z=c_{12}$ and then we compute 
\begin{align*}
    \Sigma_{a_1}=c_{12}+c_{13}+c_{14}+c_{15}+c_{16}+b_2+b_3+b_4+b_5+b_6\\
    \Sigma_{b_2}=a_1+a_3+a_4+a_5+a_6+c_{12}+c_{23}+c_{24}+c_{25}+c_{26}, \mbox{ and }\\
    \Sigma_{c_{12}}=a_1+a_2+b_1+b_2+c_{34}+c_{35}+c_{36}+c_{45}+c_{46}+c_{56},
\end{align*}
from which (\ref{eq:tritangent1}) follows. From this, we immediately obtain that $f^*(\OO_{\PP^1}(1))\cong \vartheta^{3}$.
\end{proof}

\subsection{$E_6$-covers of genus $10$ as a boundary divisor in $\overline{\cA}_6$}\label{subsect_a6}

The $E_6$-covers of genus $10$ considered in this paper fill up a boundary divisor in the compactification of the moduli space $\overline{\cA}_6$ of principally polarized abelian varieties of dimension $6$. Indeed, the main result of \cite{ADFIO} established the existence of a dominant Prym-Tyurin map $PT\colon \overline{\cH}_{E_6, 24}\dashrightarrow \overline{\cA}_6$, from the moduli space of admissible $E_6$-covers of $\PP^1$ branched over $24$ points. It turns out that pairs of  genus $10$ covers of $E_6$-type fill-up naturally the boundary divisor $E_{12:1^{27}}\subseteq \overline{\cH}_{E_6, 24}$. Its general point corresponds to an admissible $E_6$-cover
$$\bigl[ f\colon C_1\cup C_2\longrightarrow R_1\cup_t R_2, \  \ t_1+\cdots+t_{12}, \ \  s_1+\cdots+s_{12}\bigr].$$
Here $C_1$ and $C_2$ are smooth curves of genus $10$ and $f_{|C_1}\colon C_1\rightarrow R_1$, respectively $f_{|C_2}\colon C_2\rightarrow R_2$ are $E_6$-covers branched over $t_1+\cdots+t_{12}$ (respectively over $s_1+\cdots+s_{12}$). Over $t$ there is no ramification and the curves $C_1$ and $C_2$ meet in $27$ points. The arithmetic genus of the curve $C_1\cup C_2$ equals $g(C_1)+g(C_2)+27-1=46$. The Prym-Tyurin degeneration corresponding to such a degenerate admissible cover is computed in \cite[Corollary 7.5]{ADFIO} and we obtain that the corresponding semi-abelian variety $PT(f)\in \overline{\cA}_6$ has toric rank $6$. This is line with Theorem \ref{prop:corresp1}.

Pairs of $E_6$-curves depend on $2\cdot \mbox{dim}(\cH_{E_6})=18$ moduli. Adding two more dimensions for choosing the point $t$ both on $R_1$ and on $R_2$, we obtain, as expected $20=\mbox{dim}(\cA_6)-1$ moduli. 




\subsection{The $E_6$-curve of double sixes.}
Up to conjugation, the group $W(E_6)$ has a unique maximal subgroup $G_{27}$ of index $27$ corresponding to the stabilizer of a line on the reference cubic surface $X$. To fix notation, we set $G_{27}:=\mbox{Stab}_{W(E_6)}(a_6)$. Similarly, $W(E_6)$ has up to a conjugation a unique maximal subgroup $G_{36}$ of index $36$, which can be thought of as the stablizer of a double six on $X$.  Given a curve $[f\colon C\rightarrow \PP^1]\in \cH_{E_6}$, we denote by $\widetilde{f}\colon \widetilde{C}\rightarrow \PP^1$ its Galois closure and we let $C_{36}:=\widetilde{C}/G_{36}$. The induced cover 
\begin{equation}\label{eq:double-sixers}
f_{36}\colon C_{36}\longrightarrow \PP^1
\end{equation}
parametrizes double sixes, or equivalently, pairs of roots $\pm r$  on the fibres on original cover $f\colon C\rightarrow \PP^1$. Over each of the branch points $t_1, \ldots, t_{12}\in \PP^1$ of $f$,  ten pairs of roots get identified, thus the ramification profile is given by the partition $(2^{10}, 1^{16})$. From the Riemann-Hurwitz formula, we obtain  $g(C_{36})=25$.

\begin{proposition}\label{eq:aprodu-cond}
Each $E_6$-cover $f\colon C\rightarrow \PP^1$ satisfies $\mathrm{dim}\ W^1_6(C)=1$. Furthermore, we have a cover $C_{36}\rightarrow W^1_6(C)$.
\end{proposition}
\begin{proof}
This has been proved in the introduction, see (\ref{eq:pencils1}).    
\end{proof}


\begin{remark}\label{rmk:scorza} We obtain from Proposition \ref{prop:corresp1} that the incidence correspondence $\Sigma$ can be identified with the \emph{Scorza correspondence} $\mbox{Sc}(C,\vartheta)$ associated to the non-effective spin structure $\vartheta$, see \cite{DK}, \cite{FI}, that is,
$$\Sigma=\mathrm{Sc}(C,\vartheta):=\bigl\{(x,y)\in C\times C: h^0(C, \vartheta(x-y))\neq 0\bigr\}.$$ If $\theta_{\vartheta}$ is the section on  $\mbox{Pic}^0(C)$ whose zero locus is the translate of the symmetric theta divisor
$$\Theta_{\vartheta}:=\bigl\{\zeta\in \mbox{Pic}^0(C): h^0(C, \zeta\otimes \vartheta)\neq 0\bigr\},$$ then following \cite{FI}, the \emph{Szeg\H{o} kernel} $\mathrm{sz}_{\vartheta}$ is the function on $C\times C$
\begin{equation}\label{eq:szego}
\mathrm{sz}_{\vartheta}(x,y):=\frac{\theta_{\vartheta}(x-y)}{\theta_{\vartheta}(0)}.
\end{equation} The Scorza curve $S(C,\vartheta)$ is then  the vanishing locus of the Szeg\H{o} kernel.
\end{remark}

By tensoring the exact sequence 
$$0\longrightarrow \OO_C(-\Delta)\longrightarrow \OO_C\longrightarrow \OO_{\Delta}\longrightarrow 0$$
with $\omega_C\boxtimes \omega_C$, we obtain the identification 
$H^0\bigl(C\times C, (\omega_C\boxtimes \omega_C)(-\Delta)\bigr)^{S_2}\cong I_2(\omega_C)$, with the space of quadrics containing the canonical curve $C\subseteq \PP H^0(C,\omega_C)^{\vee}$. The following object plays an essential role in the entire construction.

\begin{proposition}\label{prop:scorza_quadric}
To an $E_6$-cover $f\colon C\rightarrow \PP^1$, one can associate a smooth quadric $Q\in I_2(\omega_C)$ whose associated bilinear form satisfies the equivalence for pairs $(x,y)\in C\times C$
$$Q(x,y)=0\Longleftrightarrow x=y, \mbox{ or } (x,y)\in S.$$
\end{proposition}
\begin{proof}
We use that $C\times_{\PP^1} C=\Sigma+S+\Delta$ and that $C\times_{\PP^1} C$ is linearly equivalent to $\vartheta^{3}\boxtimes \vartheta^{3}$.
Since $\OO_{C\times C}(\Sigma)\cong \vartheta\boxtimes \vartheta(\Delta)$, see also \cite[Proposition 7.1.5]{DK}, we have 
$$\OO_{C\times C}(\Delta+S)\cong \omega_C\boxtimes \omega_C(-\Delta).$$
We consider the section $Q\in H^0\bigl(C\times C, \omega_C\boxtimes \omega_C(-\Delta)\bigr)$ defined via (\ref{eq:Scorza_quadric}) and we conclude that $Q$ satisfies (\ref{eq:polar_quad}).

\vskip 4pt

We now show that $Q$, regarded as a symmetric bilinear form on $H^0(C, \omega_C)^{\vee}$ is of maximal rank. We fix a point $x\in C$ and the associated degree $10$ divisor $\Sigma_x$. In terms of lines on the cubic 
surface $X$, we may assume $x=a_1$, therefore 
\begin{equation}\label{eq:sx}
\Sigma_x=c_{12}+b_{12}+c_{13}+b_{13}+c_{14}+b_{14}+c_{15}+b_{15}+c_{16}+b_{16}.
\end{equation}
Observe that the lines $c_{1j}$ and $b_k$ are incident if and only if $j=k$. Setting $\lambda_{j}:=Q(c_{1j},b_j)\neq 0$ for $j=2, \ldots, 6$, the $10\times 10$ Gramm determinant of $Q$ in the basis of $H^0(C,\omega_C)^{\vee}$ given by the elements of $\Sigma_{a_1}$ in the order they appear in (\ref{eq:sx}) has then  the following form
$$
\bigoplus_{j=2}^6 \left(\begin{array}{rr}
0& \lambda_j\\
\lambda_j & 0\\
\end{array}
\right),
$$
which is clearly non-singular. This shows that $Q$ has maximal rank.
\end{proof}

To complete the proof of Theorem \ref{thm:main} we need to determine the Clifford index of an $E_6$-cover. Remarkably the answer is uniform on the entire Hurwitz space $\cH_{E_6}$:

\begin{proposition}\label{prop:cliff_index}
Every $E_6$-curve $C$ of genus $10$ has Clifford index $\mathrm{Cliff}(C)=4$.
\end{proposition}
\begin{proof}
We have already seen in Proposition \ref{prop:corresp1} that $W^1_6(C)\neq \emptyset$. We  show that $\mbox{gon}(C)=6$. Indeed, suppose $A\in W^1_d(C)$ is a base point free pencil with $d\leq 5$ and denote by $h\colon C\rightarrow \PP^1$ the induced cover. Retaining the notation from Proposition \ref{prop:corresp1}, since $H^0(C,\vartheta)=0$, clearly $H^0(C, \vartheta\otimes A^{\vee})=0$, therefore from Riemann-Rich we obtain $h^0(C, A\otimes \vartheta)=d$. We shall reach a contradiction, by constructing six independent sections in $H^0(C, A\otimes \vartheta)$. 

\vskip 4pt

For degree reasons, the residual correspondence $S\subseteq C\times C$ cannot be contained in the fibre product $C\times_h C$. In particular, we pick a point $x\in C$ such that if we label the points of $f^{-1}(f(x))$ as usual by $a_i, b_j, c_{ij}$, with $x=a_1$, then $h(a_i)\neq h(b_i)$, for $i=1, \ldots, 6$. Set $Z_i:=\Sigma_{a_i}+h^{-1}(h(a_i))-a_i\in |A\otimes \vartheta|$ and denote by $s_i\in H^0(C, A\otimes \vartheta)$ a section such that $\mbox{div}(s_i)=Z_i$, for $i=1,\ldots, 6$. Then, it is immediate to check that $s_i(q_j)=0$ for $i\neq j$, whereas $s_i(q_i)\neq 0$. This implies that the sections $s_1, \ldots, s_6$ are linearly independent, a contradiction.

\vskip 4pt

Assume there exists a line bundle $L$ computing $\mbox{Cliff}(C)\leq 3$, with $h^0(C,L)=r+1\geq 2$ and $d\leq 9$ and we choose $r$ minimal with this property.  Since $\mbox{gon}(C)=6$, it follows that $r\geq 3$ and $L$ is very ample. If $r\geq 4$, we contradict the Castelnuovo bound, therefore $h^0(C,L)=4$ and $\mbox{deg}(L)=9$. Clearly, the curve $\phi_L\colon C\hookrightarrow \PP^3$ cannot lie on a quadric, hence $h^0(C,L^2)\geq 10$. By Riemann-Roch, $h^0(C, \omega_C\otimes L^{-2})\geq 1$, hence for degree reasons $L^2=\omega_C$. Choose $A=A_{xy}\in W^1_6(C)$, as in (\ref{eq:pencils1}), where $(x,y)\in S$ correspond to a pair of skew lines in $f^{-1}(f(x))$, say $x=a_1$ and $y=b_1$. From the Base Point Free Pencil Trick, we obtain $h^0(C,L\otimes A^{\vee})=1$, that is, both divisors $a_1+\cdots+a_6$ and $b_1+\cdots+b_6$ span at most a plane in the embedding $\phi_L$. Since the curve $C_{36}$ of double sixers in the fibres of $f$ maps onto $W^1_6(C)$, by acting with the roots $\alpha_{ijk}$ it follows that also the divisors $a_i+a_j+a_k+c_{\ell m}+c_{m n}+ c_{n \ell}$, where $\{\ell, n, n\}=\{i, j,k\}^c$ also span a plane. We obtain in this way an effective divisor $D_{21}=\sum_{i=1}^6 a_i+\sum_{i<j} c_{ij}$ on $C$ such that $H^0(C,L(-D_{21}))\neq 0$, which is obviously impossible.  Therefore $\mbox{Cliff}(C)=4$. 
\end{proof}


\subsection{Special pencils of cubic surfaces and $E_6$-covers.} We now describe a parametrization of the Hurwitz space $\cH_{E_6}$ in terms of curves of lines on pencils of cubic surfaces. If we consider a  pencil spanned by two general cubic surfaces $X_0:=\{F_0(x,y,z,w)=0\}$ and $X_1:=\{F_1(x,y,z,w)=0\}$, this will contain $32$ singular fibres, that is, the curve of lines in the pencil spanned by $F_0$ and $F_1$ leads to an $E_6$-cover of $\PP^1$ branched over $32$ (rather than $12$) points. For a general pencil, the base curve is a smooth $(3,3)$ complete intersection in $\PP^3$, which has genus $10$. To remedy the discrepancy in the number of branch points, we consider pencils of cubics having a $10$-nodal base curve.

\vskip 4pt

We start with a $10$-nodal irreducible rational curve $B\subseteq \PP^3$ of degree $9$ which does not lie on any quadrics and denote by $\nu\colon \PP^1\rightarrow B$ the normalization map. We consider the multiplication map 
$$
\mu_B\colon \mbox{Sym}^3 H^0(B, \OO_B(1))\longrightarrow H^0(B, \OO_B(3)).
$$
Since $p_a(B)=10$, from Riemann-Roch  $h^0(B,\OO_B(3))=18$, hence 
$\mbox{dim}(\mbox{Ker}(\mu))\geq 2$. It is easy to see that for a general such curve $B$, the map $\mu_B$ is surjective, therefore there exists a pencil 
\begin{equation}\label{eq:lambda_B}
\Lambda_B:=\PP\bigl(H^0(\PP^3, \I_{B/\PP^3}(3))\bigr)
\end{equation}
of cubic surfaces containing $B$. In order to parametrize all such curves $B$, we set up the following incidence correspondence

\xymatrix{
  & \X:=\Bigl\{(Y,B): Y\in |\OO_{\PP^3}(3)|,  \ B\in |-K_Y|\mbox{ is }  10-{\mbox{nodal}}\Bigr\}\dblq PGL(4)
   \ar[dl]_{\pi_1} \ar[dr]^{\pi_2} & \\
   \mathfrak{M}_{\mathrm{cubic}} & & \mathfrak{M}_{\mathrm{rat}}      \\
                 }

\vskip 5pt

\noindent where $\mathfrak{M}_{\mathrm{cubic}}:=|\OO_{\PP^3}(3)|\dblq PGL(4)$ is the $4$-dimensional moduli space of cubic surfaces and $\mathfrak{M}_{\mathrm{rat}}$ is the quotient by $PGL(4)$ of the Hilbert scheme of $10$-nodal  rational curves of degree $9$ in $\PP^3$. The generic fibre of $\pi_1^{-1}([Y])$ is the Severi variety of $10$-nodal curves in the linear system $|-K_Y|$, which is known to be irreducible of dimension $8=\mbox{dim } |-K_Y|-10=8$, see \cite{Te}. In particular, $\X$ is irreducible of dimension $12$. The fibre $\pi_2^{-1}([B])$ is precisely the pencil $\Lambda_B$, therefore $\mathfrak{M}_{\mathrm{rat}}$ is irreducible of dimension $11$.

\vskip 4pt


We fix $[B]\in \mathfrak{M}_{\mathrm{rat}}$. For a point $\lambda\in \Lambda_B$  we denote by $X_{\lambda}\subseteq \PP^3$ the corresponding cubic surface. Let us consider the curve of lines on the cubic surfaces in $\Lambda_B$, that is, we let $C$ to be the normalization of the following curve
$$L:=\Bigl\{(\lambda, \ell): \lambda\in \Lambda_B, \  \ell\subseteq X_{\lambda} \  \mbox{ is a line}\Bigr\}.$$
The forgetful map $L\ni (\lambda,\ell)\mapsto \lambda \in \PP^1$ induces an $E_6$-cover $f\colon C\rightarrow \PP^1$. This assignment induces a rational map 
\begin{equation}\label{eq:Hurw_par}
\phi\colon \mathfrak{M}_{\mathrm{rat}}\dashrightarrow \cH_{E_6}, \ \ \ \ [B]\mapsto [f\colon C\longrightarrow \PP^1].
\end{equation}

\vskip 4pt


The fact that $\phi$ is well-defined follows from the following result:

\begin{theorem}\label{thm:curve_of_lines}
Suppose $[B]\in \mathfrak{M}_{\mathrm{rat}}$ is a general and let $C$ be the curve of lines associated to the pencil $\Lambda_B$. The following hold:
\begin{enumerate}
\item $C$ is a smooth curve of genus $10$ and the $E_6$-cover
$f\colon C\rightarrow \PP^1$ is branched over  $12$ points, with the local monodromy at each branch point, being given by a root of $W(E_6)$.
\item $C$ is isomorphic to the normalization of the curve $\mathrm{Tris}(B)$
of trisecant lines to $B\subseteq \PP^3$.
\end{enumerate}
\end{theorem}
\begin{proof}
We may assume that the curve $B\subseteq \PP^3$ does not lie on any quadric and choose $(F_0, F_1)$ to be general elements in $\PP(\mbox{Ker}(\mu_B))$. Then $B$ is necessarily the complete intersection of $F_0$ and $F_1$, in particular $\mbox{Ker}(\mu_B)$ is $2$-dimensional. 

\vskip 3pt

We now let $p_1, \ldots, p_{10}$ be the nodes of $B$.
Then the gradient vectors $\bigl(\nabla F_0\bigr)(p_i)$ and $\bigl(\nabla F_1\bigr)(p_i)$ are proportional (and non-zero), therefore there exists a unique value $\lambda_i\in \PP^1$ such that $\bigl(\nabla F_0)(p_i)=\lambda_i\bigl(\nabla F_1(\lambda_i)\bigr)$. It follows that the cubic surface $X_{\lambda_i}:=Z(F_0-\lambda_iF_1)$ is  singular at $p_i$, for each $i=1, \ldots, 10$. 
Observe that the pencil $\Lambda_B$ is tangent to the discriminant in $|\OO_{\PP^3}(3)|$ of singular cubic surfaces, as a result the local monodromy of the cover $f\colon C\rightarrow \PP^1$ at each of the points $\lambda_1, \ldots, \lambda_{10}\in \PP^1$ is the square of a reflection in $W(E_6)$, that is, trivial.
\vskip 3pt

To establish  (ii), we observe that if $\ell\subseteq X_t$ is a line, then since $B\in |-K_{X_t}|$, it follows that $B\cdot \ell=3$, that is, $\ell\in \mbox{Tris}(B)$. Conversely, if $\ell\subseteq \PP^3$ is a trisecant line to $B$, then the cubic polynomials $F_{0|\ell}$ and $F_{1|\ell}$ are proportional having the same roots, therefore there exists a unique value $\lambda_i\in \PP^1$ such that $F_{0|\ell}=\lambda_i\cdot F_{1|\ell}$ and then $\ell\subseteq X_{\lambda_i}:=Z(F_0-\lambda_iF_1)$. The discriminant of the pencil having degree $32$, we observe that $20=2\times 10$ values are absorbed in the intersection at $\lambda_1, \ldots, \lambda_{10}$, hence there exists $12$ remaining points where the corresponding cubic surface is singular at a point that does not correspond to $B_{\mathrm{sing}}$. It remains to argue that for a general $[B]\in \mathfrak{M}_{\mathrm{rat}}$, the pencil $\Lambda_B$ contains precisely $12$ singular members with a singularity that does not belong to $B_{\mathrm{sing}}$ and that each such cubic surface is $1$-nodal. To that end, we specialize to the case when the curve $B\in |-K_Y|$ becomes the following union of $9$ lines in the notation of \cite{Dol}
$$(a_1+b_6+c_6)+(c_{12}+a_2+b_1)+(b_2+c_{26}+a_6).$$
This curve has $9$ rational components and $18$ nodes and it is a union of three triples of lines forming a tritangent plane to $Y$, hence it has genus $10$. Smoothing the following eight nodes 
$$a_1\cdot b_6, \ \ b_6\cdot c_{16}, \ \ a_2\cdot b_1, \ \ b_2\cdot a_6, \ \  a_1\cdot c_{12}, \ \ b_6\cdot c_{26}, \ \ a_2\cdot c_{26}, \ \ b_1\cdot a_6$$
and preserving the remaining ten, produces a family of rational $10$-nodal curves $B_t\subseteq \PP^3$, for which $\Lambda_{B_t}$ verifies all desired properties.

\vskip 3pt

It follows that the curve of lines $f\colon C\rightarrow \PP^1$ associated to a general $[B]\in \mathfrak{M}_{\mathfrak{rat}}$ is ramified over $12$ distinct point, with the local monodromy being given by a root of $W(E_6)$. From the Riemann-Hurwitz formula we conclude that $2g(C)-2=6\cdot12-2\cdot 27=18$, therefore $g(C)=10$. 
\end{proof}

We now show that the $10$-nodal rational space curves of degree $9$ parametrize the $E_6$-Hurwitz space that will play a key role in disproving Green's Conjecture:

\begin{theorem}\label{thm:parametrization}
The morphism $\phi\colon \mathfrak{M}_{\mathrm{rat}}\dashrightarrow \cH_{E_6}$ is dominant. Therefore, the general $E_6$-curve of genus $10$ appears as the curve $\mathrm{Tris}(B)$ of trisecant lines to a $10$-nodal rational space curve $B\subseteq \PP^3$ of degree $9$. 
\end{theorem}
\begin{proof}
Recalling that the branch morphism $\mathfrak{br} \colon \cH_{E_6} \rightarrow \cM_{0,12}$ is finite, it will suffice to show that the composition $\widetilde{\phi}:=\mathfrak{br}\circ \phi\colon \mathfrak{M}_{\mathrm{rat}}\dashrightarrow \cM_{0,12}$ which assigns to each rational curve $B$ the discriminant divisor of the pencil $\Lambda_B$, is dominant. This is achieved by showing that the differential of $\widetilde{\phi}$ is generically surjective. We choose a basis $F_0, F_1\in H^0(\PP^3, \OO_{\PP^3}(3))$ for the pencil $\Lambda_B$. This pencil is then given by equations $X_{\lambda}=\lambda F_0-F_1\in H^0\bigl(\PP^3\times \PP^1, \OO_{\PP^3\times \PP^1}(3,1)\bigr)$, where $\lambda \in \PP^1$. The pencils $\Lambda_B$ has $10$ assigned singularities$(p_i,\lambda_i)$ for $i=1, \ldots, 10$ where $p_i\in B_{\mathrm{sing}}$ and $12$ unassigned singular points $(q_j, t_j)$, where $q_j\notin B$. It is now routine to check that if we set $Z:=(p_1,\lambda_1)+\cdots+(p_{10}, \lambda_{10})+(q_1,t_1)+\cdots+(q_{12},t_{12})$, then the vanishing 
$$
H^1\bigl(\PP^3\times \PP^1, \I_Z(3,1)\bigr)=0
$$
is a sufficient condition for the differential of $\widetilde{\phi}$ at the point $[B]$ to be surjective. Equivalently, the evaluation map
$\mathrm{ev}_Z\colon H^0(\PP^3\times P^1, \OO_{\PP^3\times\PP^1}(3,1))\longrightarrow H^0(\OO_Z(3,1))$ needs to have maximal rank $22$. This is a standard task that can be checked with Macaulay. 
\end{proof}



\section{The Scorza bundle of an $E_6$-cover.}
In this section, we introduce a quadratic bundle naturally determined by the Scorza quadric, and explain its relations with the linear syzygies of a canonical $E_6$-curve.



\subsection{The quadric bundle of an $E_6$-cover.} We regard the smooth Scorza quadric $Q\in I_2(\omega_C)$ constructed in Proposition \ref{prop:scorza_quadric} as a linear isomorphism  $Q^{\sharp}\colon H^0(C,\omega_C)^{\vee}\rightarrow H^0(C,\omega_C)$. The canonical embedding of $C$ induces for each $x\in C$ a $1$-dimensional space 
$$
\ell_x :=\omega_{C|x}^{\vee}\subseteq H^0(C,\omega_C)^{\vee}
$$ 
generated by the evaluation map $\mathrm{ev}_x\colon H^0(C,\omega_C)\rightarrow \mathbb C$. We consider the $1$-dimensional space of holomorphic forms  
\begin{equation}\label{eq:hx}
\langle h_x\rangle :=Q^{\sharp}\bigl(\omega_{C|x}^{\vee})\subseteq H^0(C,\omega_C).
\end{equation}

In local coordinates, if $(\omega_1,\ldots,  \omega_g)$ is a basis of $H^0(C,\omega_C)$ and $Q=\sum_{i,j=1}^g a_{ij}\omega_i \omega_j$, then 
$$h_x=\sum_{i,j=1}^g a_{ij}\omega_i(x)\omega_j\in H^0(C,\omega_C).$$
Using Proposition \ref{prop:corresp1}, we observe that $\mbox{div}(h_x)=S_x+2x$, for every $x\in C$. The holomorphic form $h_x$  is intimately related to the \emph{Scorza quartic} $F(C,\vartheta)\in \mbox{Sym}^4 H^0(C,\omega_C)$, considered in \cite{DK}, \cite[Theorem 6.2]{FI}. Up to a non-zero constant, this is the unique quartic in $\PP^{g-1}\cong \PP(H^0(C,\omega_C)^{\vee})$ having the remarkable property that the mixed polars $P_yP_z F(C,\vartheta)$ with respect to pairs $(y,z)\in \mbox{Sc}(C,\vartheta)$ are quadrics of rank one. One has the following precise connection between $h_x$ and the Scorza quartic, which also serves as a justification for referring to $Q$ as the Scorza quadric.

\begin{proposition}\label{prop:scorza_quar}
Keeping the notation as above, let $x\in C$ and choose points $y,z\in f^{-1}(f(x))$ such that $x,y,z$ correspond to a tritangent triangle. Then,  
one has the following identity:
$$P_yP_z F(C,\vartheta)=h_x^2\in \mathrm{Sym}^2 H^0(C,\omega_C).$$
\end{proposition}
\begin{proof}
We pick $x,y,z\in C$ distinct points in the same fibre of $f$ such that $x, y,z$ form a tritangent plane, that is,  $(x,y), (y,z), (z,x)\in \Sigma$. We write therefore
$$\Sigma_y=z+x+B_{yz}, \ \mbox{ and } \ \Sigma_z=y+x+B_{zy},$$
where $\mbox{deg}(B_{yz})=\mbox{deg}(B_{zy})=8$. It suffices to show that 
\begin{equation}\label{eq:sc2}
B_{yz}+B_{zy}=S_x,
\end{equation}
which  amounts to showing that $\mbox{div}(h_x)=\mbox{supp} (\vartheta(y-z))+\mbox{supp}(\vartheta(y-z))$. This follows from (\ref{eq:tritangent1}), since $f^{-1}(f(x))=x+\Sigma_x+S_x$.
\end{proof}



\vskip 4pt

We introduce the dual quadric $q:=Q^{\vee}\in \mbox{Sym}^2 H^0(C,\omega_C)^{\vee}$, which can be regarded as a bilinear form $q\colon H^0(C,\omega_C)\otimes H^0(C,\omega_C)\rightarrow \mathbb C$ given by the inverse matrix $(a_{ij})_{i,j=1, \ldots, g}^{-1}$. Clearly, 
$$q\bigl(h_x, H^0(C, \omega_C(-x)\bigr)=0,$$
 and, in fact, we have 
\begin{equation}\label{eq:hxperp}
    H^0(C, \omega_C(-x)\bigr)= \langle h_x\rangle^\perp,
\end{equation}
where orthogonality is defined with respect to the quadratic form $q$. 

%Geometrically, $\mathbf P (H^0(C,\omega_C(-x)))$ is the tangent hyperplane at the point $[h_x]\in \mathbf P(H^0(C,\omega_C))$ to the quadric defined by $q$ in $\mathbf P(H^0(C,\omega_C))$.} 


\begin{definition}\label{def:quadric_bundle}
 The \emph{Scorza bundle} on $C$ is
defined as  $\cH_Q:=\Omega^1_Q|_C\otimes \omega_C$. 
\end{definition}
One has an exact sequence of vector bundles on $C$
\begin{equation}\label{eq:quadr_bun1}
0\longrightarrow \omega_C^{\vee}\longrightarrow M_{\omega_C}\longrightarrow \cH_Q\longrightarrow 0.
\end{equation}

For $x\in C$, we have therefore the identification of the fibre $\cH_{Q|x}\cong  H^0\bigl(C, \omega_C(-x)\bigr)/\langle h_x\rangle$. 
\begin{definition}\label{def:v_y}
For  $y\in S_x$, we have that $h_y\in H^0(C, \omega_C(-x))$  by Proposition \ref{prop:scorza_quadric} and \eqref{eq:hxperp},
and we set 
$$v_y:=h_y \  \mathrm{mod }\  \langle h_x\rangle \in \cH_{Q|x}.$$
\end{definition}


The  Scorza bundle $\cH_Q$ is equipped with a symmetric bilinear form, which slightly abusing notation, we also denote by $$q\colon \cH_Q\otimes \cH_Q\longrightarrow \OO_C,$$ which fibrewise is the bilinear map $q_x\colon \cH_{Q|x}\otimes \cH_{Q|x}\rightarrow \mathbb C$ induced by the bilinear form which we also denoted by $q\colon H^0(C,\omega_C)\otimes H^0(C,\omega_C)\rightarrow \mathbb C$. In this way, the Scorza bundle becomes a quadratic bundle on $C$, and $\mbox{det}(\mathcal H_Q)=\mathcal O_C$.



\subsection{The Scorza bundle and $5$th linear syzygies.}
The bilinear form $q$ on $\cH_Q$ induces a bilinear form  $q_k\colon \bigwedge^k \cH_Q \otimes \bigwedge^k \cH_Q\rightarrow \OO_C$, by setting 
$$q_k\bigl(u_1\wedge\cdots\wedge u_k, v_1\wedge\cdots\wedge v_k\bigr):=\det(q(u_i,v_j))_{i,j=1, \ldots, k}.
$$

For a normalized volume element $\operatorname{vol}_{\cH_Q}\colon \bigwedge^8\cH_Q \stackrel{\sim}{\longrightarrow} \mathcal O_C$,
the  Hodge star operator 
\[
\star\colon \bigwedge^k \cH_Q\longrightarrow \bigwedge^{8-k} \cH_Q
\]
is characterized by the property
\begin{equation}\label{eq:star-definition}
% \alpha\wedge \bigl(\star\beta\bigr)=\langle\alpha,\beta\rangle\operatorname{vol}_{\cH_Q}.
  \operatorname{vol}_{\cH_Q}\big(\alpha\wedge \bigl(\star\beta\bigr)\big)=q_k(\alpha,\beta).
\end{equation}
  

In middle degree $4$, we have $\star^2=\mathrm{id}$, therefore $\bigwedge^4 \cH_Q$ splits into two eigenbundles $\mathcal F^\pm$, both  of rank $35$, corresponding to eigenvalues $\pm 1$. We next observe that any non--zero global section of one of the vector bundles $\F^\pm$, leads to a non-trivial syzygy in the range relevant for the Green Conjecture. 

\begin{proposition}
\label{prop:star}
One has a natural injection $H^0(C,\mathcal F^\pm)\lhook\joinrel\longrightarrow K_{5,1}(C,\omega_C)$.
\end{proposition}
\begin{proof}
By taking exterior powers in (\ref{eq:quadr_bun1}), then tensoring with $\omega_C$, we obtain an inclusion
\begin{equation}\label{eq:middle-inclusion}
 \bigwedge^4\cH_Q\lhook\joinrel\longrightarrow\bigwedge^5 M_{\omega_C}\otimes \omega_C,
\end{equation}
inducing an injection $
H^0\bigl(C,\mathcal \bigwedge^4 \cH_Q\bigr)\lhook\joinrel\longrightarrow H^0\bigl(C, \bigwedge^5 M_{\omega_C}\otimes \omega_C\bigr)$.  Note that we have the inclusion $H^0\bigl(C,\bigwedge^5 M_{\omega_C}\otimes \omega_C)\subseteq \bigwedge^5 H^0(C,\omega_C)\otimes H^0(C,\omega_C)$.

\vskip 3pt

We now show that $H^0(C,\F^{\pm})\subseteq H^0\bigl(C,\bigwedge^4 \cH_Q\bigr)$ intersects trivially $d_{6,0}\bigl(\bigwedge^6H^0(C,\omega_C)\bigr)$.
To that end, assume $d_{6,0}(\tau)=\gamma\in H^0(C, \F^{\pm})$, with $\tau\in\bigwedge^6H^0(C,\omega_C)$.  For an isotropic evaluation functional $x\in H^0(C,\omega_C)^\vee$, the boundary has value $\iota_x\tau$.  A hyperbolic decomposition $H^0(C,\omega_C)=\mathbb C u\oplus\mathbb C v\perp H_0$ with $u=Q^\sharp x$ and $q(u,v)=1$, identifies the space $H_0$ with $\mathcal H_{Q|x}$ and the map in~\eqref{eq:middle-inclusion} with $\eta\mapsto u\wedge\eta$.  The two Hodge star maps satisfy the following identity, for a certain scalar $\epsilon$:
\begin{equation}\label{eq:null-Hodge}
 \star_{H^0(\omega_C)}(u\wedge\eta)=\varepsilon\,u\wedge\star_{H_0}\eta.
\end{equation}

Both Hodge-star operators, as well as the inclusion \eqref{eq:middle-inclusion} vary algebraically with $x\in C$, therefore so does $\varepsilon$. Since $\varepsilon^2=-1$, we obtain that $\varepsilon$ is constant.

 Let $\delta\in\{1,-1\}$ be the fixed sign such that
$\gamma\in H^0(C,\mathcal F^\delta)$.
Under the preceding identifications, for each $x\in C$ we have
\[
\iota_{\ell_x}\tau=u\wedge\eta,
\qquad
\star_{H_0}\eta=\delta\eta,
\]
where $u=Q^\sharp(\ell_x)$ and $\eta$ represents
$\gamma|_x$. Applying~\eqref{eq:null-Hodge}, with $\star$ denoting
the Hodge operator on
$\bigwedge^5H^0(C,\omega_C)$, we obtain
$$
\star(\iota_{\ell_x}\tau)=\star(u\wedge\eta)
=\varepsilon\,u\wedge\star_{H_0}\eta
=\varepsilon\delta\,u\wedge\eta
=\varepsilon\delta\,\iota_{\ell_x}\tau.
$$
Since both $\varepsilon$ and $\delta$ are
independent of $x$, setting
$\lambda:=\varepsilon\delta$ gives
\[
\star(\iota_{\ell_x}\tau)
=\lambda\,\iota_{\ell_x}\tau,
\qquad\text{for every }x\in C.
\]





Since the canonical curve $C$ spans $\PP H^0(C,\omega_C)^\vee$, it follows that $\star(\iota_{\ell}\tau)=\lambda \iota_{\ell}(\tau)$, for each   $\ell\in H^0(C,\omega_C)^\vee$.  Contracting by $\ell$, yields
\[
 Q^\sharp(\ell)\wedge\iota_{\ell}\tau=0\in \bigwedge^6 H^0(C,\omega_C), \ \mbox{ for each  } \ \ell\in H^0(C,\omega_C)^\vee.
\]
This implies $\tau=0$. Indeed, we choose an orthonormal basis $\omega_1,\ldots,\omega_{10}$ of $H^0(C, \omega_C)$ and set $\ell_j:=q(\omega_j,-)\in H^0(C,\omega_C)^{\vee}$.  Since $\tau\in \bigwedge^6 H^0(C,\omega_C)$, summing the equations above,
\[
 0=\sum_{j=1}^{10}\omega_j\wedge\iota_{\ell_j}\tau=6\tau,
\]
forcing $\tau=0$. This finishes the proof. 
\end{proof}




\begin{corollary}\label{cor:neq 0}
    Suppose that $H^0\bigl(C,\bigwedge^4\cH_Q\bigr)\ne 0$. Then $K_{5,1}(C,\omega_C)\ne 0$.
\end{corollary}



The quadratic form $q$ on $\mathcal H_Q$ induces
a nondegenerate quadratic form $q_2$ on
$\bigwedge^2\mathcal H_Q$. Its inverse defines
a nonzero global section
$
q_2^{-1}\in H^0\bigl(C,
\operatorname{Sym}^2(\bigwedge^2\mathcal H_Q)\bigr).
$
However, its image under exterior multiplication
$
\operatorname{Sym}^2(\bigwedge^2\mathcal H_Q)
\longrightarrow \bigwedge^4\mathcal H_Q
$
vanishes identically. Thus the quadratic structure alone does not produce the nonzero four-form required by
Corollary~\ref{cor:neq 0}. The sixteen isotropic planes
constructed above allow us to remedy this, by considering a distinguished subspace of alternating forms on which the induced quadratic
form has a nondegenerate restriction. The inverse of this restricted form will give rise to the desired four-form. Accordingly, we introduce a new sheaf $\mathcal N$ next that will model the family of distinguished subspaces of
alternating forms.









\section{$D_5$-isotropic planes of the Scorza bundle}

An essential reason for the presence of extra syzygies on $E_6$-curves is the presence of a $W(D_5)$-structure on the Scorza bundle $\cH_Q$. Recall that a del Pezzo surface  $Y$ of degree $4$ can be obtained from our reference cubic surface $X$ by contracting one line. To fix notation, we choose to contract the line $b_2$ and the map giving this contraction is given by the linear system $$|5h-2a_1-a_2-2a_3-2a_4-2a_5-2a_6|$$ on $X$. Denoting by $S_{\mathrm{dP}_4}\subseteq \PP^4$ the image of $X$, ten of the lines of $X$ are mapped to the lines of $S_{\mathrm{dP}_4}$, which are indexed by $a_2, b_1,b_3,b_4,b_5, b_6$ and $c_{ij}$, for $3\leq i<j\leq 6$. The incidence of lines on $S_{\mathrm{dP}_4}$ is the same as on $X$, in particular every line of $Y$ meets five lines. The root lattice $D_5:=K_{S_{\mathrm{dP}_4}}^{\perp}$ has $40$ roots. It is well-known that $S_{\mathrm{dP}_4}$ has two conic bundle structures $S_{\mathrm{dP}_4}\rightarrow \PP^1$, each having the following configuration of degenerate fibres
\begin{equation}\label{eq;dp4}
\bigl(b_3+c_{13}, b_4+c_{14}, b_5+c_{1}, b_6+c_{16}\bigr)\  \mbox{ and } \  \bigl(b_1+a_2, c_{34}+c_{56}, c_{35}+c_{46}, c_{36}+c_{45}\bigr).
\end{equation}

\subsection{The sixteen isotropic planes and the bundle
of alternating forms}

We retain the notation of previous section and recall that for $x\in C$, we have  $\cH_{Q|x}=H^0(C,\omega_C(-x))/\langle h_x\rangle$ and for $y\in S_x$ we have defined the vectors $v_y\in \cH_{Q|x}$ via Definition \ref{def:v_y}. Recall  the notation (\ref{eq:pencils1})
$$\Sigma_y=D_{xy}+R_{xy}'\in C_6,$$
where $R_{xy}'$ denote the divisor on $C$ corresponding to the $5$ lines that are incident to $y$ and skew to $x$. For pairs $(x,y)\in S$ such that $f(x)=f(y)$ is not a branch point, we set 
$$
L_{xy}:=
\langle v_z:z\in R_{xy}'\rangle
\subseteq \mathcal H_{Q|x}.
$$

\begin{proposition}\label{prop:l_xy}
$L_{xy}$ is a $4$-dimensional isotropic subspace of $\cH_{Q|x}$.
\end{proposition}
\begin{proof}
To simplify notation, we may assume $x=b_2$ and $y=b_1$, in which case $$R_{b_2 b_1}'=a_2+c_{13}+c_{14}+c_{15}+c_{16}\in C_5.$$ Observe that the five lines in $R_{b_2 b_1}'$ are skew, therefore $q(L_{b_2 b_1}, L_{b_2 b_1})=0$. Furthermore, observe that $\OO_C(b_2+R_{b_2 b_1}')\in W^1_6(C)$, therefore 
$\mbox{dim } \langle b_2, R_{b_2 b_1}'\rangle =4$, from which it follows that the vectors $v_{a_2}, v_{c_{13}}, v_{c_{14}}, v_{c_{15}}, v_{c_{16}}$ span a $4$-dimensional subspace of $\cH_{Q|b_2}$. 
\end{proof}

We now set
\begin{equation}\label{eq:p_xy}
P_{xy}:=
\langle v_y\rangle+
\bigl(L_{xy}\cap \langle v_y\rangle ^{\perp}\bigr)\
\subseteq\mathcal H_{Q|x}.
\end{equation}

\begin{proposition}
$P_{xy}$ is also a $4$-dimensional isotropic subspace of $\cH_{Q|x}$.      
\end{proposition}
\begin{proof} It suffices to observe that the $L_{xy}\ni v\mapsto q(v_y,v)$ is a non-zero  map, hence $L_{xy}\cap \langle v_y\rangle ^{\perp}$ is $3$-dimensional.
\end{proof}

The description (\ref{eq:p_xy}) of the $4$-planes $P_{xy}$ has to be extended to the case when the points $x, y$ lie over one of the branch points $t_1, \ldots, t_{12}$ of $f$. Recall that  $\pi_1\colon S\to C$ is the first projection of the skew correspondence. Set $S^{\circ}:=S\setminus (f\circ \pi_1)^{-1}\bigl(\{t_1, \ldots, t_{12}\bigr\}\bigr)$. Since the relative orthogonal Grassmannian bundle $OG(4, \pi_1^*\cH_Q)$ is proper,
the map $S^{\circ}\rightarrow OG(4, \pi_1^*\cH_Q)$ given by  $(x,y)\mapsto P_{xy}$ extends  to define a maximal isotropic rank $4$ subbundle
\begin{equation}\label{eq:P4}
\mathcal P\subseteq \pi_1^*(\mathcal H_Q).
\end{equation}




\vskip 4pt


We introduce the sheaf
\begin{equation}\label{eq:N}
\mathcal N:=
\ker\left\{
\bigwedge\nolimits^2\mathcal H_Q^\vee
\longrightarrow
(\pi_1)_*\bigwedge\nolimits^2\mathcal P^\vee
\right\}.
\end{equation}
Thus $\mathcal N_{|x}$ consists of the alternating
forms vanishing on all $16$ isotropic $4$-planes
$P_{xy}$, where $y\in S_x$. We shall show that $\cN$ is locally free of rank $3$. Due to the rigid structure of the planes $P_{xy}$, there is therefore a three dimensional family of skew-symmetric forms on $\cH_{Q|x}$ vanishing on all these planes. Before proving this we need some preparation. Our presentation will be very explicit and uses only facts on the Picard lattice of del Pezzo surfaces, although a more conceptual representation-theoretic proof is certainly possible.

\vskip 4pt

We fix $x\in C$ such that $f(x)$ is not a branch point. We choose a pair of distinct points $y,z\in S_x$, such that $(y,z)\in \Sigma$, that is, the lines labeled by $y$ and $z$ meet. Since the action of $W(E_6)$ on such triples is transitive, we may assume $x=b_2, y=b_1$ and $z=a_2$ and we set 
\begin{equation}\label{eq:E-F}
E:=P_{b_2 b_1} \ \mbox{ and } \ \ F:=P_{b_2 a_2}.
\end{equation}
Since these are isotropic $4$-spaces with respect to $q_x=q$ such that $E\cap F=\{0\}$, we have
\begin{equation}\label{eq:H-splitting}
    \cH_{Q|b_2}\cong E\oplus F\cong E\oplus E^{\vee}.
\end{equation}
We choose a hyperbolic basis of $\cH_{Q|b_2}$ with respect to this decomposition as follows. Set 
\begin{equation}\label{eq:e_i,f_i}
e_i:=v_{c_{1 i}}, \ \   f_i:=v_{b_i} \in \cH_{Q|b_2}, \  \mbox{ for } \ \ i=3, \ldots, 6.     
\end{equation}
Note that we can rescale these vectors in such a way that $q(e_i, e_j)=q(f_i, f_j)=0$ and $q(e_i, f_j)=\delta_{ij}$. Furthermore, we can also rescale $v_{b_1}$ such that
$$v_{b_1}=f_3+f_4+f_5+f_6\in \cH_{Q|b_2} \ \ \mbox{ and } \ \ v_{a_2}=r_3e_3+r_4e_4+r_5e_5+r_6e_6\in \cH_{Q|b_2},$$ where $r_i=q(v_{a_2}, f_i)\neq 0$. Note that since $b_1\cdot a_2=1$, it follows that $q(v_{b_1}, v_{a_2})=r_3+\cdots+r_6\neq 0$.

\begin{lemma}\label{lemma:E+F}
With this notation, we have the following equations for the isotropic planes $E$ and $F$:
$$E=\bigl\langle v_{b_1}, e_3-e_6, e_4-e_6, e_5-e_6\bigr\rangle \ \mbox{ and } \ \ F=\bigl\langle v_{a_2}, r_4f_3-r_3f_4, r_5f_3-r_3f_5, r_6f_3-r_3f_6 \bigr\rangle.
$$
\end{lemma}

We will express all $16$ isotropic spaces $P_{b_2 y}$, for $y\in S_{b_2}$ in terms of this basis. For $\{i,j,k\ell\}=\{3,4,5,6\}$, we denote by $\epsilon_{ijk\ell}$ the signature of the permutation 
$$
\begin{pmatrix}
3 & 4 & 5 & 6 \\
i & j & k & \ell
\end{pmatrix}
$$
We have the following:

\begin{lemma}\label{lemma:16iso}
For different integers $i,j\in \{3,4,5,6\}$, we set $\{k,\ell\}=\{i,j\}^c$. Similarly, for a given $i$, we set $\{j,k, \ell\}=\{i\}^c$. We have the following description for the $14$ remaining isotropic planes $P_{b_2 y}\subseteq \cH_{Q|b_2}$:
\begin{enumerate}
\item $P_{b_2 b_i}=\bigl\langle f_i, -e_j+\epsilon_{ijk\ell}(r_kf_{\ell}-r_{\ell} f_k), -e_k+\epsilon_{ijk\ell}(r_{\ell} f_j-r_j f_{\ell}), -e_{\ell}+\epsilon_{ijk\ell}(r_jf_k-r_kf_j)\bigr\rangle$.
\item $P_{b_2 c_{1i}}=\bigl\langle e_i, \ f_j+f_k+f_{\ell}, \ e_j-e_k+\epsilon_{ijk\ell} r_i f_{\ell}, \ e_j-e_{\ell}-\epsilon_{ijk\ell}r_if_k\bigr\rangle$.
\item $P_{b_2 c_{ij}}=\bigl \langle e_i-e_j-\epsilon_{ijk\ell} r_{\ell}f_k+\epsilon_{ijk\ell} r_kf_i, \ f_i+f_j, \ e_k+\epsilon_{ijk\ell}r_{\ell}f_i, \ e_{\ell}-\epsilon_{ijk\ell} r_kf_i\bigr\rangle$.
\end{enumerate}
\end{lemma}
\begin{proof}
This is a direct computation, using the $W(E_6)$-incidence of the lines on the cubic surface $X$, translated into the property that $q(v_y, v_z)=0$, unless $y$ and $z$ are incident lines.    
\end{proof}

We now describe the space of alternating forms $\alpha \in \bigwedge^2 \cH_{Q|b_2}^{\vee}$ vanishing of the $16$ isotropic planes $P_{b_2 y}$. We will single out six of these planes, namely
\begin{equation}\label{eq:six_planes}
E=P_{b_2b_1}, \ F=P_{b_2a_2}, \ P_{b_2 b_3}, \ P_{b_2b_4}, \ P_{b_2b_5}, \ P_{b_2b_6}.
\end{equation}
We show that a skew-form vanishing on these six planes actually vanishes on all $16$ planes. To that end, we introduce the space of skew-forms vanishing on the planes in (\ref{eq:six_planes}:


\begin{definition}\label{def:N6}
Set $N_6:=\Bigl\{\alpha\in \bigwedge^2 \cH_{Q|b_2}^{\vee}: \alpha_{|F\times F}=0, \  \ \alpha_{|P_{b_2b_i}\times P_{b_2b_i}}=0, \ \mbox{ for } i=1, 3, \ldots, 6\Bigr\}.$
\end{definition}

We have the following description of this space, in terms of the coordinates 

\begin{equation}\label{eq:coordinates}
A_{ij}=\alpha(e_i, e_j), \ C_{ij}=\alpha(f_i,f_j) \mbox{ and } B_{ij}=\alpha(e_i, f_j),  \mbox{ for } i, j=3, \ldots, 6,
\end{equation}
with respect to the chosen hyperbolic basis of the space $\cH_{Q|b_2}$:

\begin{lemma}\label{lemma:N6}
The space $N_6$ is three-dimensional. For any skew-form $\alpha\in N_6$, there exist complex numbers $(m_3, m_4, m_5, m_6)$ unique up to addition of one common constant, such that in terms of the coordinates (\ref{eq:coordinates}), one has
\begin{align}\label{eq:A_ij,B_ij}
A_{ij}=m_i-m_j,     \ C_{ij}=\frac{m_i-m_j}{r_k r_{\ell}}, \ B_{ij}=-\frac{\epsilon_{ijk\ell}}{r_i}\bigl(m_k-m_{\ell}\bigr), \ \mbox{ for } i\neq j, \mbox{ and } \ B_{ii}=0.
\end{align}
\end{lemma}
\begin{proof}
We begin by recalling that $E$ is generated by the vectors $v_{b_1}, e_3-e_6, e_4-e_6$ and $e_5-e_6$. Writing the condition $\alpha(e_i-e_6, e_j-e_6)=0$, we obtain $A_{ij}-A_{i6}+A_{j6}=0$, therefore there exist $m_3, \ldots, m_6\in \mathbb C$ such that $A_{ij}=m_i-m_j$. Of course, the $m_i$'s are determined up to addition by a common constant.

\vskip 3pt
From the relation $\alpha(e_i-e_6, v_{b_1})=0$, we obtain that 
$$\sum_{j=3}^6 B_{3j}=\sum_{j=3}^6 B_{4j}=\sum_{j=1}^5 B_{5j}=\sum_{j=3}^6 B_{6j}.$$

We now exploit that $\alpha_{|F\times F}=0$ and recall that the basis of $F$ is described by Lemma \ref{lemma:E+F}. We obtain the following relations
$$r_i C_{jk}=r_j C_{ik}-r_k C_{ij}, \mbox{ for distinct } i, j, k\in \{3,4,5,6,\}.$$
These relations can be rewritten as saying that there exist complex numbers $n_3, \ldots, n_6$ with 
\begin{equation}\label{eq:Cij}
C_{ij}=r_i n_j-r_j n_i, \mbox{ for } 3\leq i<j\leq 6.
\end{equation}

We now use that $\alpha_{|P_{b_2 b_i}\times P_{b_2 b_i}}=0$. From the relation $\alpha\bigl(f_i, -e_j+\epsilon_{ijk\ell}(r_kf_{\ell}-r_{\ell}f_k)\bigr)=0$, we obtain that 
\begin{equation}\label{eq:Bij}
B_{ji}=\epsilon_{ijk\ell} r_i\bigl(r_{\ell} n_k-r_k n_{\ell}\bigr),
\end{equation}
that is, all entries $B_{ji}$ away from the diagonal are determined by the quantities $n_3, \ldots, n_6$. Furthermore, all diagonal entries $B_{33}, \ldots, B_{66}$ are equal and we denote their common value by $B$. Then we use that $\alpha\bigl(-e_j+\epsilon_{ijk\ell}(r_k f_{\ell}-r_{\ell} f_k), -e_k+\epsilon_{ijk\ell}(r_{\ell} f_j-r_jf_{\ell})\bigr)=0$, for all quadruples 
$i, j, k, \ell$, to obtain the relations 
$r_i r_j r_{\ell} n_j-r_i r_j r_{\ell} n_k=2\epsilon_{ijk\ell} r_{\ell}B$. Writing this relation with $i$ and $\ell$ interchanged, but keeping $j$ and $k$ fixed, we obtain that $(r_i+r_{\ell})B=0$. Since this holds for every pairs $i\neq \ell$, in particular we obtain $(r_3+r_4+r_5+r_6)B=0$. Since $r_3+r_4+r_5+r_6\neq 0$, we conclude $B=0$, that $B_{ii}=0$. Furthermore, we also note that the quantity $r_j r_k r_{\ell} n_i-m_i$ is independent of $i$. Since the quantities $n_i$ admit a change of variables of the type $n_i\mapsto n_i+tr_i$, for $i=3, \ldots, 6$, we obtain that we may choose $n_i=\frac{m_i}{r_j r_k r_{\ell}}$. Substituting, into (\ref{eq:Bij}) and into (\ref{eq:Cij}), we finish the proof. 
\end{proof}

We show that the $3$-dimensional space $N_6$ defined by vanishing along $6$ isotropic planes, actually characterizes vanishing of skew-symmetric forms on $\cH_{Q|b_2}$ along all the $16$ isotropic hyperplanes induced by the $W(D_5)$-structure.

\begin{proposition}\label{prop:6=16}
We have the following equality of $3$-dimensional subspaces in $\bigwedge^2 \cH_{Q_2|b_2}$:
$$N_6=\Bigl\{\alpha\in \bigwedge^2\cH_{Q|b_2}^{\vee}: \alpha_{|P_{b_2 y}\times P_{b_2,y}}=0, \mbox{ for every } y\in S_{b_2}\Bigr\}.$$
\end{proposition}
\begin{proof}
 We take a skew-symmetric form $\alpha\colon \cH_{Q|b_2}\times \cH_{Q|b_2}\rightarrow \mathbb C$ having coefficients $A_{ij}$, $B_{ij}$ and $C_{ij}$ as in Lemma \ref{lemma:N6}. We directly verify that the remaining then isotropic planes $P_{b_2, c_{1j}}$ and $P_{b_2, c_{ij}}$ impose no further linear conditions. We shall carry this out only for the 
 plane $P_{b_2, c_{13}}$, which using Lemma \ref{lemma:16iso}, is generated by the vectors $e_3$, $f_4+f_5+f_6$, $e_4-e_5+r_3f_6$ and $e_4-e_6-r_3f_5$. We check directly that the form $\alpha$ vanishes on these vectors, for instance 
 $$\alpha\bigl(e_3, f_4+f_5+f_6\bigr)=B_{34}+B_{35}+B_{36}=0 \mbox{ and } \alpha\bigl(e_3, e_4-e_5+r_3f_6\bigr)=A_{34}-A_{35}+r_3 B_{36}=0.$$
 The remaining verifications are similar.
\end{proof}

We now return to the sheaf $\cN$ defined via (\ref{eq:N}). We have the following:
\begin{corollary}  
\label{cor:rank-three}
The sheaf $\mathcal N$ is a rank-three locally free subbundle of
$\bigwedge^2\mathcal H_Q^\vee$. 
\end{corollary}
\begin{proof}
By definition, $\cN$ is a torsion free subsheaf of the vector bundle $\bigwedge^2 \cH_Q^{\vee}$ and the cokernel is also torsion free. Hence both $\mathcal N$ and the cokernel are locally free. Proposition \ref{prop:6=16} shows that the rank of $\mathcal N$ equals three.     
\end{proof}





Lemma \ref{lemma:N6} shows that one has an isomorphism, 
\begin{equation}\label{eq;N6_isom}
\mathbf C^4/\bigl\langle (1,1,1,1)\bigr \rangle \xrightarrow{\ \sim\ }N_6 .
\end{equation}
For $m=(m_3,\dots,m_6)\in\mathbf C^4$ we denote by $\alpha_m$ the alternating form on $\mathcal \cH_{Q|b_2}$ whose coordinates \eqref{eq:coordinates} are given by \eqref{eq:A_ij,B_ij}. 
For $i=3,\dots,6$, let $\alpha^{(i)}:=\alpha_m$, where $m$ is the $i$-th unit vector. Since $\sum_{i=3}^6 \alpha^{(i)}=0$, it follows that  
$\alpha^{(3)},\alpha^{(4)},\alpha^{(5)}$ is a basis of $N_6$.

\vskip 4pt

\begin{remark}\label{rmk:d5}
Although we have not spelled this out, the incidence of the $16$ isotropic planes $P_{b_2 y}$, where $y\in S_{b_2}$ is governed by the two conic bundle structures (\ref{eq;dp4}) on the del Pezzo surface $S_{\mathrm{dP}_4}$. Indeed, one has a relation $(b_3-c_{13})+(b_4-c_{14})+(b_5-c_{15})+(b_6-c_{16})=2(b_1-a_2)$ in $\mbox{Pic}(S_{\mathrm{dP}_4})$, therefore 
$$\bigl\langle b_3-c_{13}, b_4-c_{14}, b_5-c_{15}, b_6-c_{16}\bigr\rangle/\bigl\langle b_1-a_2\bigr\rangle \stackrel{\cong}\longrightarrow \mathbb C^4/\bigl\langle (1,1,1,1)\bigr\rangle \cong N_6,$$
which uses the description (\ref{eq;N6_isom}) provided by Proposition \ref{prop:6=16}, which accordingly could be alternatively established via representation-theoretic means.
\end{remark}

The next proposition completes the analysis over points of $x\in C$, not lying over a branch point of $f$, and provides the last ingredient needed in the construction of non-zero linear syzygies.

\begin{proposition}\label{prop:q2formula}
For $m=(m_3, m_4, m_5, m_6)$, we have the formula
$$
q_2\bigl(\alpha_m, \alpha_m\bigr)=\frac{4}{r_3r_4r_5r_6}\sum_{3\le i<j\le 6}r_ir_j\,(m_i-m_j)^2.
$$
In particular $q_2$ restricts to a nondegenerate form on $N_6=\cN_{|b_2}$. The Gram determinant with respect to $\alpha_{e_3},\alpha_{e_4},\alpha_{e_5}$ equals $64(r_3+r_4+r_5+r_6)^2/(r_3r_4r_5r_6)^2\neq 0$.
\end{proposition}

\begin{proof}
We start by describing $q_2$ in hyperbolic coordinates. Consider the basis $(e^i,f^i)$ of $\cH_{Q|b_2}^\vee$ dual to $(e_i,f_i)$, given by $e^i=q(f_i, -)$ and $f^i=q(e_i, -)$. Up to the normalizations of \eqref{eq:e_i,f_i}, these are the classes of the evaluation functionals at $b_i$ and at $c_{1i}$, respectively. In particular this basis is hyperbolic for $Q$:
\[
Q(e^i,f^j)=\delta_{ij},\qquad Q(e^i,e^j)=Q(f^i,f^j)=0.
\]
An alternating form with coordinates \eqref{eq:coordinates} is given by 
\[
\alpha=\sum_{i<j}A_{ij}\,e^i\wedge e^j+\sum_{i<j}C_{ij}\,f^i\wedge f^j+\sum_{i,j}B_{ij}\,e^i\wedge f^j\in \bigwedge^2 \cH_{Q|b_2}^{\vee}.
\]
By the determinantal definition of $q_2$, taken with respect to $Q$,
\[
q_2\bigl(e^i\wedge e^j,f^k\wedge f^\ell\bigr)=\delta_{ik}\delta_{j\ell}-\delta_{i\ell}\delta_{jk},\qquad q_2\bigl(e^i\wedge f^j,e^k\wedge f^\ell\bigr)=-\delta_{i\ell}\delta_{jk},
\]
and all other pairings of these basis vectors vanish. Hence
\begin{equation}\label{eq:q2hyp}
q_2(\alpha,\alpha)=2\sum_{3\le i<j\le 6}A_{ij}C_{ij}-\sum_{i,j=3}^{6}B_{ij}B_{ji}.
\end{equation}

Recall that  for $i\neq j$, we denote $\{k,\ell\}:=\{i,j\}^c$. By Lemma \ref{lemma:N6}, $A_{ij}C_{ij}=(m_i-m_j)^2/(r_kr_\ell)$. We also have $B_{ij}B_{ji}=-(m_k-m_\ell)^2/(r_ir_j)$. Substituting into \eqref{eq:q2hyp} gives
\[
q_2(\alpha_m,\alpha_m)=2\sum_{i<j}\frac{(m_i-m_j)^2}{r_kr_\ell}+2\sum_{i<j}\frac{(m_k-m_\ell)^2}{r_ir_j}.
\]
The two sums coincide after replacing each pair by its complement, which yields the desired expression. To prove non-degeneracy, note that the bilinear form associated with $q_2|_{N_6}$ is equal to $4\Phi/(r_3r_4r_5r_6)$, where
\[
\Phi(m,m')=\bigl(r_3+\cdots+r_6\bigr)\sum_{i=3}^6 r_im_im'_i-\Bigl(\sum_{i=3}^6 r_im_i\Bigr)\Bigl(\sum_{i=3}^6 r_im'_i\Bigr).
\]
The conclusion now follows from a routine calculation.
%If $m$ lies in the radical of $\Phi$, then $u\,r_im_i=r_i\sum_jr_jm_j$ for every $i$. Since $r_i\neq0$ and $u\neq0$, this forces $m_i=\sum_jr_jm_j/u$ for all $i$, i.e.\ $m$ is constant and $\alpha_m=0$.
%For the determinant, observe that in the basis $\alpha_{e_3},\alpha_{e_4},\alpha_{e_5}$ the Gram matrix of $\Phi$ is $u\operatorname{diag}(r_3,r_4,r_5)-(r_3,r_4,r_5)^t(r_3,r_4,r_5)$. By the matrix determinant lemma its determinant is $u^2r_3r_4r_5r_6$. Multiplying by $\bigl(4/(r_3r_4r_5r_6)\bigr)^3$ gives the stated value.
\end{proof}





\subsection{The $16$ isotropic planes over ramification and antiramification points.}\label{antiram}

The $E_6$-cover $f\colon C\rightarrow \PP^1$ is branched over $12$ points $t_1, \ldots, t_{12}$ and we explain the structure and the incidence of the $16$ planes $P_{xy}$, when $f(x)=t_1$ is a branch point and $y\in S_x$. Assume, without loss of generality, that the root describing the local monodromy around $t_1$ is $\alpha_{\max}$, therefore the sheets of $f$ that come together are those labeled by the pairs $(a_1, b_1)$, \ldots, $(a_6,b_6)$. The remaining $15$ points in the fibre $f^{-1}(t_i)$ are then labeled by $c_{ij}$, with $i\neq j$, and we refer to these points as being the \emph{anti-ramification} points of $f$. If $x=a_1=b_1$ is a ramification point, then it is easy to see that $\pi_1\colon S\rightarrow C$ is unramified over $f(a_1)=f(b_1)$. For an antiramification point, say $x=c_{12}$, since $a_3, \ldots, a_6, b_3, \ldots, b_6\in S_{c_{12}}$, we obtain that the fibre $\pi_1^{-1}(c_{12})$ corresponds to the partition $(4^2, 1^{8})$.

\vskip 4pt

If $x=b_2=a_2$ is a ramification point, then in the limit we still have $16$ limiting isotropic planes $P_{b_2 y}$. A simple calculation shows that the limiting hyperplane $P_{a_2b_2}$ is now given by 
$$P_{b_2a_2}=\bigl\langle r_3e_3+r_4 f_4+r_5f_5+r_6e_6,  r_4f_3-r_3f_4,  r_5f_3-r_3f_5,  r_6f_3-r_3f_6\bigr\rangle.$$
The definition of the hyperbolic basis $e_i, f_j$ remains the same. Similarly the formulas for the remaining $14$ planes $P_{b_2 y}$ are those given in Lemma \ref{lemma:16iso}. If $x=c_{12}$ is an anti-ramification point, then some of the isotropic planes coalesce. First, we observe that the isotropic planes $P_{c_{12} a_i}$ and $P_{c_{12} c_{1i}}$ all coalesce to the plane $F=\langle f_3, \ldots, f_6\rangle$, but the remaining $8$ planes stay distinct. So in this case, we have a configuration of $9$ (rather than $16$) distinct isotropic planes.


The analysis of the limiting incidence planes over the branch values of $f$ leads to the same conclusion: the restriction of $q_2$ to 
$\mathcal{N}_{|x}$ is non-degenerate at every point $x\in C$. At ramification points this follows verbatim from the arguments above. At the antiramification points, the computation is analogous. We omit the details.



















\section{A global $4$-form on the Scorza bundle}

The Scorza bundle $\cH_Q$ is equipped with the quadratic form $q\colon \cH_Q\otimes \cH_Q\rightarrow \OO_C$. Recall that we have constructed the rank $3$ subvector bundle $\cN\subseteq \bigwedge^2 \cH_Q^{\vee}$. In particular, every local section $\alpha$ of $\cN$ induces a skew symmetric local map $A_{\alpha}\colon \cH\rightarrow \cH$.  This assignment leads to a symmetric form
$$\mathfrak{b}\colon \cN\otimes \cN\longrightarrow \OO_C, \ \ \ (\alpha, \beta)\mapsto \mathrm{tr}(A_{\alpha} A_{\beta}).$$
In particular, $\mathfrak{b}$ induces a morphism $\mathfrak{b} \colon \cN\rightarrow \cN^{\vee}$.




\begin{lemma}\label{lem:nondegeneracy-incidence-net}
The morphism $\mathfrak{b} \colon\mathcal N\to\mathcal N^\vee$ is an isomorphism. In particular, $\mathfrak{b}^{-1}$ is a nowhere vanishing section of $\operatorname{Sym}^2\mathcal N$.
\end{lemma}
\begin{proof}
If $f(x)$ is not a branch value, then $\cN_{|x}=N_6$, and $\mathfrak{b}_{|x}$ is non-degenerate by Proposition~\ref{prop:q2formula}. Over the branch values, $\mathfrak{b}_{|x}$ is non-degenerate by the discussion in \ref{antiram}. Thus $\mathfrak{b}$ is an isomorphism on every fibre.
\end{proof}







We now consider the composition 
$$\cN\otimes \cN\hookrightarrow \bigwedge^2 \cH_{Q}^{\vee}\otimes \bigwedge^2 \cH_Q^{\vee}\longrightarrow \bigwedge^4 \cH_Q^{\vee}.$$
Since this map is obviously symmetric, it induces a morphism $\iota\colon \mbox{Sym}^2 \cN\rightarrow \bigwedge^4 \cH_Q^{\vee}$. Note that $\bigwedge^4 \cH_Q^{\vee}\cong \bigwedge^4 \cH_Q$.

\begin{definition}\label{def:tautsect}
The  \emph{Scorza section} is defined as $\Omega:=H^0(\iota)(\mathfrak{b}^{-1})\in H^0\bigl(C,\bigwedge^4 \cH_Q\bigr)$.
\end{definition}




\begin{proposition}\label{prop:scorza-section}
The Scorza section $\Omega$ is non-zero. In particular, $K_{5,1}(C,\omega_C)\ne 0$.
\end{proposition}

\begin{proof}
By Lemma \ref{lem:nondegeneracy-incidence-net}, $\mathfrak{b}^{-1}\in H^0\bigl(C,\operatorname{Sym}^2\mathcal N\bigr)$. Hence $\widetilde\Omega:=\iota(\mathfrak{b}^{-1})$ is a global section of $\bigwedge^4\mathcal H_Q^\vee$, and $\Omega$ is then its image under the isomorphism $\bigwedge^4\mathcal H_Q^\vee\cong\bigwedge^4\mathcal H_Q$ induced by $q$. It therefore suffices to show that $\widetilde\Omega$ does not vanish at one point. We take $x=b_2$ with $f(x)$ not being a branch point, and use the notation of the previous section.

We first compute $b^{-1}$ at $x$. By Proposition~\ref{prop:q2formula}, for every $m\in\mathbb C^4$ and $i=3,\dots,6$,
\[
\mathfrak{b}\bigl(\alpha_m,\alpha^{(i)}\bigr)=\frac{4}{r_3r_4r_5r_6}\,r_i\Bigl((r_3+\cdots+r_6)\,m_i-\sum_{j=3}^6 r_jm_j\Bigr).
\]
Since $\alpha_m=\sum_{i=3}^6 m_i\alpha^{(i)}$ and   $\alpha^{(3)}+\alpha^{(4)}+\alpha^{(5)}+\alpha^{(6)}=0$, it follows that
\[
\frac{r_3r_4r_5r_6}{4(r_3+r_4+r_5+r_6)}\sum_{i=3}^6\frac{1}{r_i}\,\mathfrak{b}\bigl(\alpha_m,\alpha^{(i)}\bigr)\,\alpha^{(i)}=\sum_{i=3}^6\Bigl(m_i-\frac{1}{(r_3+\cdots+r_6)}\sum_{j=3}^6 r_jm_j\Bigr)\alpha^{(i)}=\alpha_m .
\]
Hence
\[
\mathfrak{b}^{-1}_{|x}=\frac{r_3r_4r_5r_6}{4(r_3+r_4+r_5+r_6)}\sum_{i=3}^{6}\frac1{r_i}\,\alpha^{(i)}\otimes\alpha^{(i)}
\quad\text{and}\quad
\widetilde\Omega_{|x}=\frac{r_3r_4r_5r_6}{4(r_3+r_4+r_5+r_6)}\sum_{i=3}^{6}\frac1{r_i}\,\alpha^{(i)}\wedge\alpha^{(i)}.
\]

We now evaluate on $(e_3,f_3,f_4,f_5)$. For an alternating form $\alpha$ with coordinates \eqref{eq:coordinates}, we have $\alpha(e_3,f_3)=B_{33}=0$. With the usual shuffle convention for the wedge product this gives
\[
(\alpha\wedge\alpha)\bigl(e_3,f_3,f_4,f_5\bigr)=2\bigl(B_{35}C_{34}-B_{34}C_{35}\bigr).
\]
For $\alpha=\alpha_m$, Lemma \ref{lemma:N6} yields
\begin{multline*}
(\alpha_m\wedge\alpha_m)(e_3,f_3,f_4,f_5)\\
=-2\Bigl(\frac{(m_4-m_6)(m_3-m_4)}{r_3r_5r_6}+\frac{(m_5-m_6)(m_3-m_5)}{r_3r_4r_6}\Bigr).
\end{multline*}
This expression vanishes for $\alpha^{(3)}$ and $\alpha^{(6)}$. It equals $2/(r_3r_5r_6)$ for $\alpha^{(4)}$ and $2/(r_3r_4r_6)$ for $\alpha^{(5)}$. Therefore
\[
\widetilde\Omega_{|x}(e_3,f_3,f_4,f_5)=\frac{1}{r_3r_4r_5r_6}\neq0 .
\]
Thus $\Omega\neq0$, and Corollary \ref{cor:neq 0} implies that $K_{5,1}(C,\omega_C)\neq0$.
\end{proof}




















\section{A counterexample to the Green-Lazarsfeld Secant Conjecture}

We finish the paper with the announcement of a counterexample to the Green-Lazarsfeld Secant Conjecture \cite{GL} in genus $10$, for the $E_6$-curves treated throughout the paper. We have the following result:

\begin{theorem}\label{thm:gl}
Let $[f\colon C\rightarrow \PP^1]\in \cH_{E_6}$ be a smooth $E_6$-cover of genus $10$. Then there exist $4$-very ample line bundles $L\in\mathrm{Pic}^{20}(C)$ such that $K_{3,2}(C,L)\neq 0$.
\end{theorem}

In other words, we construct a line bundle $L\in \mbox{Pic}^{20}(C)$ such that $L-\omega_C\notin C_5-C_3$, with $K_{3,2}(C,L)\neq 0$. From the point of view of the difference variety appearing in the conjectured equivalence (\ref{conj:GL}), this is therefore a codimension $2$ counterexample of the Secant Conjecture. Note moreover that in this case, since we have established in Proposition \ref{prop:cliff_index} that $\mbox{Cliff}(C)=4$, we have an equality $d=2g+p+1-\mbox{Cliff}(C)$. Full details of the proof of Theorem \ref{thm:gl} will be presented shortly in a separate paper, here we discuss the construction.

\vskip 4pt

We pick $t_0\in \PP^1$ and label the points of $f^{-1}(t_0)$ according to the notation from \cite{Dol}. In particular, $(a_1, \ldots, a_6)$ and $(b_1, \ldots, b_6)$ form a double sixer corresponding to the root $\alpha_{\mathrm{max}}$. Set 
$$M:=\vartheta(a_1+\cdots+a_6)\in \mbox{Pic}^{15}(C).$$\begin{proposition}\label{prop:M}
We have $H^1(C,M)=0$ and $h^0(C,M)=6$. Moreover $M$ is globally generated but not very ample and the induced $\phi_M\colon C\rightarrow \PP^5$ identifies the pairs $(a_1, b_1), \ldots, (a_6, b_6)$.
\end{proposition}

The image $D:=\phi_M(C)$ is a $6$-nodal curve of arithmetic genus $16$ and degree $15$. We consider the multiplication map 
$$
\mu_M \colon \mbox{Sym}^2 H^0(C,M)\longrightarrow H^0\bigl(C,M^{2}\bigr).
$$
It turns out that $\mbox{dim } \mbox{Ker}(\mu_M)=5$ and the base locus of the quadrics in $\mathrm{bs}\bigl(\mbox{Ker}(\mu_M)\bigr)$ is a smooth \emph{quintic del Pezzo surface} $S_{\mathrm{dP}_5}\subseteq \PP^5$. Moreover, since $\omega_D\cong \OO_D(2)$, it follows that the $6$-nodal curve $D$ is a cubic section of $S_{\mathrm{dP}_5}$.

\vskip 3pt

Each quintic del Pezzo surface can be uniquely expressed as a linear section of the Grassmannian $G(2,V)$, where $V$ is a $5$-dimensional vector space. Precisely, we have 
$$S_{\mathrm{dP}_5}=G(2,V)\cap \PP^5\subseteq \PP\bigl(\bigwedge^2 V\bigr),$$
where the intersection is taken with respect to the Pl\"ucker embedding. This induces a stable vector bundle of rank $2$ on $C$, by taking 
$$F:=\phi_M^*\bigl(\cU_{|S_{\mathrm{dP}_5}}^{\vee}\bigr),$$
where $\cU\subseteq V\otimes \OO_{G(2,V)}$ is the tautological rank $2$ bundle. It is easy to show that $F$ is stable and 
$$\mbox{det}(F)\cong M, \ \ \ h^0(C,M)=6.$$

\vskip 4pt

We have already established that the pencil $\OO_C(a_1+\cdots+a_6)\in W^1_6(C)$ belongs to a one-dimensional of $W^1_6(C)$ parametrized by the curve $C_{36}$ of double sixes on the fibres of $f$. In particular, the Petri map for each such pencil has a $1$-dimensional kernel, that is, 
$$h^0\bigl(C, \omega_C(-2a_1-\ldots-2a_6)\bigr)=1.$$
We let $p\in \mbox{supp}\bigl(\omega_C(-2a_1-\ldots-2a_6)\bigr)$ and we perform a general positive elementary transformation of $F$ at the point $p$, that is, we have an exact sequence
$$0\longrightarrow F\longrightarrow F_1\longrightarrow \mathbb C(p)\longrightarrow 0.$$
It is easy to show that $F_1$ is a stable vector bundle with $\mbox{det}(F_1)\cong M(p)$ and $h^0(C,F_1)=5$. Thus $F_1$ belongs to the Brill-Noether locus $\mathcal{BN}_C(2, 16, 5)$ of stable rank $2$ vector bundles on $C$ having degree $16$ and at least $5$ sections. Observe that $H^0(C,F)\cong H^0(C, F_1)$, in particular $F_1$ is not globally generated at $p$. 

\vskip 3pt

We consider the line bundle $L:=\omega_C^2\otimes \bigl(\mbox{det}(F)\bigr)^{\vee}=\vartheta^3(-a_1-\cdots-a_6-p)\in \mbox{Pic}^{20}(C)$. It is easy to show that $L$ is not $4$--very ample, in particular the non-vanishing $K_{3,2}(C,L)\neq 0$, follows directly from \cite{GL}. We establish the following result:

\begin{theorem}\label{th:def_GL}
There exists a deformation $F_{\mathrm{def}}$ of the vector bundle $F_1$ inside the space $\mathcal{BN}_C(2, 16,5)$, such that  the line bundle $L_{\mathrm{def}}=\omega_C^{2}\otimes \bigl(\mathrm{det}(F_{\mathrm{def}})\bigr)^{\vee}\in \mathrm{Pic}^{20}(C)$ satisfies
$$L_{\mathrm{def}}-\omega_C\notin C_5-C_3, \ \mbox{  and } K_{3,2}(C,L_{\mathrm{def}})\neq 0.$$
\end{theorem}

\begin{remark}
The construction presents an interesting connection to \cite[Theorem 1.7]{FK}, where it is shown that on a Petri general curve $C$ of genus $g$ any  counterexample $L\in \mbox{Pic}^d(C)$ to the Secant Conjecture in the case $d=2g+p+1-\mbox{Cliff}(C)$, necessarily leads to the secant variety 
$$V_{g-p-3}^{g-p-4}\bigl(\omega_C^2\otimes L^{\vee}\bigr)$$ having dimension larger than expected. In our case $\omega_C^2\otimes L^{\vee}=\vartheta(a_1+\cdots+a_6+p)=M(p)$ and, indeed, the variety of secants 
$$V_4^3\bigl(M(p)\bigr)=\bigl\{D\in C_4: h^0(C, M(p-D))\geq 4\bigr\} $$
has dimension one, since $C+a_i+b_i+p\subseteq V_4^3\bigl(M(p)\bigr)$, for $i=1, \ldots, 6$. However this condition is lost in the deformation $L_{\mathrm{def}}$ of $L$. On the other hand, the $E_6$-curves are not Petri general, in particular they satisfy $\mbox{dim } W^1_6(C) = 1$, so the conditions of \cite[Theorem 1.7]{FK} are not satisfied.
\end{remark}



\vskip 9pt

\noindent
\emph{AI usage statement}.
ChatGPT-6 Astra was used for exploratory computations in the search for counterexamples, reference search, and interactive feedback on the manuscript. The paper has been completely written by the authors who take full responsibility for all content.




\begin{thebibliography}{aaaaaaa}

\bibitem[ACGH]{ACGH}
  E.~Arbarello, M.~Cornalba, P.\,A.~Griffiths and J.~Harris, {\emph{Geometry of algebraic curves, Vol.~I}}, Grundlehren math.\ Wiss.\ vol.~267, Springer-Verlag, New York, 1985.
\bibitem[ADFIO]{ADFIO}   V. Alexeev, R. Donagi, G. Farkas, E. Izadi, and A. Ortega,
  \emph{The uniformization of the moduli space of principally polarized abelian
  6-folds}, J. Reine Angew. Mathematik \textbf{761} (2020), 163--217.
\bibitem[Ag]{Ag} D. Agostini, {\em{The Martens-Mumford Theorem and the Green-Lazarsfeld Secant Conjecture}},  Journal of Algebraic Geometry  \textbf{33} (2024), 629--654.  
\bibitem[Ap]{Ap} M. Aprodu, {\em{Remarks on syzygies of $d$-gonal curves}}, Mathematical Research Letters \textbf{12} (2005), 387--400.
\bibitem[AF1]{AF1} M. Aprodu and G. Farkas, {\emph{The Green Conjecture for smooth curves lying on arbitrary $K3$ surfaces}}, Compositio Mathematica \textbf{147} (2011), 839--851.
\bibitem[AF2]{AF2} M. Aprodu and G. Farkas, {\emph{Koszul cohomology and applications to moduli}}, in: Grassmanians, Moduli Spaces and Vector Bundles (E. Previato ed.), Clay Mathematics Proceedings, Vol. 14, 2011, 25-50.
\bibitem[AF3]{AF3} M. Aprodu and G. Farkas, {\emph{Green's Conjecture for general covers}}, in: Compact Moduli Spaces and Vector Bundles, Contemporary Math. Vol. 564, 2012, 211--226.  
\bibitem[AFPRW]{AFPRW} M. Aprodu, G. Farkas, \c S. Papadima, C. Raicu,
and J. Weyman, {\em Koszul modules and Green's Conjecture},
Invent. Math. \textbf{218} (2019), 657--720.
\bibitem[Dol]{Dol}
I. Dolgachev, \emph{Classical algebraic geometry: A modern view}, Cambridge University
  Press, Cambridge, 2012.
\bibitem[DK]{DK} I. Dolgachev and V. Kanev, {\em{Polar covariants of plane cubics and quartics}}, Advances in Math. \textbf{98} (1993), 216-301.  

\bibitem[DS]{DS}
R. Donagi and R.  Smith, \emph{The structure of the {P}rym map}, Acta
  Math. \textbf{146} (1981), 25--102.

\bibitem[Don]{Don} R. Donagi, {\em{The unirationality of $\mathcal{A}_5$}},
Annals of Math. \textbf{119} (1984), 269--307.
\bibitem[Eis]{Eis} D. Eisenbud, {\em{The geometry of syzygies}}, Graduate Texts in Mathematics \textbf{229}, Springer 2005.
\bibitem[F]{F} G. Farkas, {\em{Difference varieties and the Green-Lazarsfeld Secant Conjecture}},  \'{E}pijournal de G\'{e}om\'{e}trie Alg\'{e}brique 2024 (special volume in honor of Claire Voisin), 11658.  
\bibitem[FI]{FI} G. Farkas and E. Izadi, {\em{Szeg\H{o} kernels and Scorza quartics on the moduli space of spin curves}}, arXiv:2409.13303, to appear in the American J. Math.
\bibitem[FK]{FK} G. Farkas and M. Kemeny, {\em The generic Green–Lazarsfeld Secant Conjecture}, Invent. Math. \textbf{203} (2016), 265–301.
\bibitem[Gr]{Gr} M. Green, {\em{Koszul cohomology and the cohomology of projective varieties}}, Journal of Diff. Geometry \textbf{19} (1984), 125--171.
\bibitem[GL1]{GL} M. Green and R. Lazarsfeld, {\em{On the projective normality of complete linear series on an algebraic curve}}, Inventiones Math. \textbf{83} (1986), 73--90.
\bibitem[GL2]{GL2} M. Green and R. Lazarsfeld, {\em{Some results on the syzygies of finite sets and algebraic curves}}, Compositio Math. \textbf{67} (1988), 301-314.
\bibitem[HR]{HR} A. Hirschowitz and S. Ramanan, {\em{New evidence for Green's Conjecture on syzygies of canonical curves}}, Annales Scientifiques de l'\'Ecole Normale Sup\'erieure \textbf{31} (1998), 145--152.
\bibitem[K1]{K1}
V. Kanev, \emph{Spectral curves, simple {L}ie algebras, and {P}rym-{T}jurin
  varieties}, Theta functions---{B}owdoin 1987, Proc. Sympos. Pure Math., vol.~49, Amer. Math. Soc.
  1989, 627--645.
\bibitem[K2]{K2}
V. Kanev, \emph{Hurwitz spaces of {G}alois coverings of {${\mathbb P}^1$}, whose
  {G}alois groups are {W}eyl groups}, J. Algebra \textbf{305} (2006),
  442--456. 
\bibitem[K]{K} M. Kemeny, {\em{Universal Secant Bundles and Syzygies of Canonical Curves}}, Inventiones Math. \textbf{223} (2021), 995--1026.
\bibitem[Mo]{Mo} F. Moretti, {\em{Brill-Noether theory and Green's conjecture for general curves on simple abelian surfaces}}, International Math. Research Notices, \textbf{7} (2025), 1--23.
\bibitem[Pa]{Pa} J. Park, {\em{Syzygies of tangent developables surfaces and $K3$ carpets via secant varieties}}, Algebra \& Number Theory \textbf{19} (2025), 1029--1047.
\bibitem[RS]{RS} C. Raicu and S. Sam, {\em{Bi-graded Koszul modules, $K3$ carpets, and Green's conjecture}}, Compositio Math. \textbf{158} (2022), 33--56.
\bibitem[Sag]{Sag} M. Sagraloff, {\em{Special linear series and syzygies of canonical curves of genus $9$}}, arXiv:0605758.
\bibitem[Sch]{Sch} F.-O. Schreyer, {\em{Syzygies of canonical curves and special linear series}}, Mathematische Annalen \textbf{275} (1986), 105-137.
\bibitem[Te]{Te} D. Testa, {\em{The irreducibility ofthe space of rational curves on del Pezzo surfaces}}, Journal of Algebraic Geometry \textbf{18} (2009), 37--61.
\bibitem[V1]{V1} C. Voisin, {\em{Green's generic syzygy conjecture
for curves of even genus lying on a $K3$ surface}}, Journal of
European Mathematical Society \textbf{4} (2002), 363--404.
\bibitem[V2]{V2} C. Voisin, {\em{Green's canonical syzygy conjecture for generic curves of odd genus}},
Compositio Mathematica \textbf{141} (2005), 1163--1190.

\end{thebibliography}


\end{document} 
